%!TEX TS-program = xelatex
%!TEX TS-options = -shell-escape
%!TEX root = ../OpAlg2_SS16/operatoralgebren2.tex
\RequirePackage{fix-cm} 
\documentclass[a4paper, twoside, headsepline, index=totoc,toc=listof,toc=bibliography,toc=index, fontsize=10pt, cleardoublepage=empty, headinclude, DIV=12, BCOR=5mm, titlepage,draft]{scrartcl}
%!TEX root = ../AnaTopGeo_SS14/ana_top_geo.tex
\usepackage{scrtime} % KOMA, Uhrzeit ermoeglicht

%--Pakete zum "Programmieren"
% ======================================================================================
\usepackage{etoolbox}
\usepackage{letltxmacro}
\usepackage{ifthen}
% ======================================================================================

%--Farbdefinitionen und Grafiken (muss vor tikz geladen werden)
% ======================================================================================
\usepackage[usenames, table, x11names]{xcolor}
\definecolor{dark_gray}{gray}{0.45}
\definecolor{light_gray}{gray}{0.6}
\definecolor{fb10_blue}{cmyk}{0.8,0.4,0.13,0.07}
\usepackage[final]{graphicx}
\usepackage{adjustbox}
\newcommand{\cfbox}[2]{% coloured frame box
	\ifmmode
	\mathchoice{\adjustbox{cfbox=#1}{$\displaystyle#2$}}{\adjustbox{cfbox=#1}{$\textstyle#2$}}{\adjustbox{cfbox=#1}{$\scriptstyle#2$}}{\adjustbox{cfbox=#1}{$\scriptscriptstyle#2$}}
	\else
	\adjustbox{cfbox=#1}{#2}
	\fi
}
% ======================================================================================

%--Zum Zeichnen/ TikZ-Kram (vor polyglossia bzw. babel geladen werden)
% ======================================================================================
\usepackage{tikz}
\usepackage{tikz-cd}
\usetikzlibrary{external}
\tikzset{>=latex}
\usetikzlibrary{%
	shapes,
	arrows.meta,
	intersections,
	calc,
	3d,
	decorations.pathreplacing,decorations.markings,decorations.pathmorphing,
	angles,
	quotes,
}
\tikzexternalize[prefix=tikz/,up to date check=diff]
\pgfkeys{/pgf/images/include external/.code=\includegraphics{#1}}
\tikzset{external/system call={lualatex \tikzexternalcheckshellescape -halt-on-error -interaction=batchmode --shell-escape -jobname "\image" "\texsource"}}
\AtBeginEnvironment{tikzcd}{\tikzexternaldisable} % tikzexternalize fuer tikzcd deaktivieren, da inkompatibel
\AtEndEnvironment{tikzcd}{\tikzexternalenable}
\tikzset{% um Inkompatibilitaeten von quotes und polyglossia bzw. babel zu vermeiden
  every picture/.append style={
    execute at begin picture={\shorthandoff{"}},
    execute at end picture={\shorthandon{"}}
  }
}
\usepackage{pgfplots}
\usepgfplotslibrary{colormaps}
\newcommand*\circled[1]{\tikzexternaldisable\tikz[baseline=(char.base)]{\node[shape=circle,draw,inner sep=2pt] (char) {#1};}\tikzexternalenable}
% ======================================================================================



%-- Mathepakete etc.
% ======================================================================================
\usepackage[T1]{fontenc}
\renewcommand{\rmdefault}{zpltlf}
\usepackage{mathtools} % beinhaltet amsmath
\mathtoolsset{showonlyrefs,centercolon,showmanualtags}
\newtagform{brackets}[\textbf]{[}{]}
\usetagform{brackets}
\usepackage{fix-cm}
\usepackage[bbgreekl]{mathbbol}
\usepackage{amssymb,marvosym} 
\usepackage{nicefrac} % schräge Brüche
\usepackage{faktor}
\newcommand{\Faktor}[1]{\faktor[\textstyle]{#1}}
\usepackage{xfrac}
\usepackage{cancel}
\usepackage{mathdots} % Verbesserung von Punkten wie zB \ldots
\usepackage[bb=px]{mathalfa} % \mathbb als px font
\usepackage{centernot}
\usepackage{stackrel}
\DeclareSymbolFont{bbold}{U}{bbold}{m}{n}
\DeclareSymbolFontAlphabet{\mathbbold}{bbold}
\newcommand{\ind}{\mathbbold{1}} % charakteristische-Funktion-Eins
\def\mathul#1#2{\color{#1}\underline{{\color{black}#2}}\color{black}} %farbiges Untersteichen im Mathe-Modus
\renewcommand{\le}{\leqslant}
\renewcommand{\ge}{\geqslant}
% ======================================================================================


%-- Von xfrac erzeuge font warnings ignorieren
% ======================================================================================
\usepackage{silence}
\WarningFilter{latexfont}{Size substitutions with differences}
\WarningFilter{latexfont}{Font shape `U/bbold/m/n' in size}
% ======================================================================================


%-- Typographie/Polyglossia
% ======================================================================================
\usepackage[euler-digits]{eulervm} % vor fontspec laden!
\usepackage[no-math]{fontspec}
\usepackage{polyglossia} % moderner babel-ersatz
\setmainlanguage[spelling=new,babelshorthands=true]{german}
\shorthandoff{"}
\setotherlanguage{english}
\defaultfontfeatures{Mapping=tex-text, WordSpace={1.2}, Ligatures={Required,Common,Contextual},Extension=.otf} %


\setmainfont{TeXGyrePagellaX}[UprightFont=*-Regular,BoldFont=*-Bold,ItalicFont=*-Italic,BoldItalicFont=*-BoldItalic,ItalicFeatures={Style=Historic},Ligatures={Required,Common,Contextual,Historic}]
\setsansfont{texgyreadventor}[Scale=MatchUppercase, UprightFont=*-regular, BoldFont=*-bold, ItalicFont=*-italic, BoldItalicFont=*-bolditalic]
\setmonofont{SourceCodePro}[Scale=0.9,UprightFont=*-Regular, BoldFont=*-Semibold, ItalicFont=*-Light]
\usepackage{xltxtra}
\usepackage{fontawesome}
\usepackage[final]{microtype}
\usepackage[draft=false]{scrlayer-scrpage} 
\flushbottom
% ======================================================================================


%-- Aufzählungen
% ======================================================================================
\usepackage[shortlabels,inline]{enumitem}
\setlist[itemize,1]{label=\faCaretRight}
\setlist[enumerate]{font=\bfseries}
\setlist[description]{font=\normalfont\bfseries}
\usepackage{multicol}
% ======================================================================================


%-- Floats/Figures/Tabellen
% ======================================================================================
\usepackage{wrapfig}
\usepackage{float}
\usepackage[margin=10pt, font=small, labelfont={sf, bf}, format=plain, indention=1em]{caption}
\captionsetup[wrapfigure]{name=Abb. }
\usepackage{booktabs}
% ======================================================================================


%-- korrekte Anführungszeichen und Zitierbefehle
% ======================================================================================
\usepackage[autostyle,german=quotes,english=british]{csquotes}
% ======================================================================================


%--Indexverarbeitung
% ======================================================================================
\usepackage{makeidx}
\newcommand{\bet}[1]{\textbf{\emph{#1}}}
\newcommand{\Index}[1]{\bet{#1}\index{#1}}
\makeindex
\setindexpreamble{{\noindent\sffamily\small Die \emph{Seitenzahlen} sind mit Hyperlinks versehen und somit anklickbar} \par \bigskip}
\renewcommand{\indexpagestyle}{scrheadings}
% ======================================================================================


%-- Marginnotes/Todonotes/Footnotes
% ======================================================================================
\deffootnote[1.5em]{1.5em}{1.5em}{\textsuperscript{\thefootnotemark}\ }
\usepackage[fulladjust]{marginnote}
\renewcommand*{\marginfont}{\itshape\footnotesize}
\usepackage[textsize=small]{todonotes}
\usepackage{ragged2e}
\renewcommand*{\raggedleftmarginnote}{\RaggedLeft}
\renewcommand*{\raggedrightmarginnote}{\RaggedRight}
\LetLtxMacro{\oldtodo}{\todo}
\renewcommand{\todo}[2][]{\tikzexternaldisable\oldtodo[#1]{#2}\tikzexternalenable}
\LetLtxMacro{\oldmissingfigure}{\missingfigure}
\renewcommand{\missingfigure}[2][]{\tikzexternaldisable\oldmissingfigure[{#1}]{#2}\tikzexternalenable}
% ======================================================================================


% -- BibLaTeX
% ======================================================================================
\usepackage[%
	backend=biber,
	sortlocale=auto,
	natbib,
	hyperref,
	backref,
	style=alphabetic
	]%
{biblatex}
\renewcommand*{\mkbibnamelast}[1]{%
  \ifmknamesc{\textsc{#1}}{#1}}
\renewcommand*{\mkbibnameprefix}[1]{%
  \ifboolexpr{ test {\ifmknamesc} and test {\ifuseprefix} }
    {\textsc{#1}}
    {#1}}
\def\ifmknamesc{%
  \ifboolexpr{ test {\ifcurrentname{labelname}}
               or test {\ifcurrentname{author}}
               or ( test {\ifnameundef{author}} and test {\ifcurrentname{editor}} ) }}
\addbibresource{../!config/quellen.bib}
% ======================================================================================

%--Konfiguration von Hyperref und Cleveref
% ======================================================================================
\usepackage[hidelinks, pdfpagelabels,  bookmarksopen=true, bookmarksnumbered=true, linkcolor=black, urlcolor=SkyBlue2, plainpages=false,pagebackref, citecolor=black, hypertexnames=true, pdfauthor={Jannes Bantje}, pdfborderstyle={/S/U}, linkbordercolor=SkyBlue2, colorlinks=false,final,backref=false]{hyperref}
\usepackage[nameinlink,noabbrev]{cleveref}
\newcommand{\appendLink}[1]{#1\,\faExternalLink}
\newcommand{\hrefsym}[2]{\href{#1}{\texttt{\appendLink{#2}}}}
\newcommand{\hrefsymX}[2]{\href{#1}{\appendLink{#2}}}
\newcommand{\hrefsymmail}[2]{\href{#1}{\texttt{\faEnvelopeO\,#2}}}
\renewcommand{\url}[1]{\hrefsym{#1}{\nolinkurl{#1}}}
% ======================================================================================


% -- QR-Codes (hinter hyperref laden!)
% ======================================================================================
\usepackage{qrcode}
% ======================================================================================

%--Römische Zahlen
% ======================================================================================
\newcommand{\RM}[1]{\MakeUppercase{\romannumeral #1{}}}
% ======================================================================================

%-- Definition von diversen Mathe-Befehlen
% ======================================================================================
%!TEX root = mitschrift_main.tex

% -- Zum Finetuning von Befehlen
% ======================================================================================
\makeatletter
\newcommand{\raisemath}[1]{\mathpalette{\raisem@th{#1}}}
\newcommand{\raisem@th}[3]{\raisebox{#1}{$#2#3$}}
\makeatother
\makeatletter
\newcommand{\killDescendersM}[1]{\mathpalette{\killD@scendersM{#1}}}
\newcommand{\killD@scendersM}[2]{\raisebox{0pt}[\height][0pt]{$#2#1$}}
\makeatother
\DeclareRobustCommand{\minwidthbox}[2]{%
  \ifmmode
    \expandafter\mathmakebox
  \else
    \expandafter\makebox
  \fi
  [\ifdim#2<\width\width\else#2\fi]{#1}%
}
% ======================================================================================


%-- Klammerbefehle
% ======================================================================================
\DeclarePairedDelimiter{\abs}{\lvert}{\rvert}
\DeclarePairedDelimiter{\floor}{\lfloor}{\rfloor}
\DeclarePairedDelimiter{\ceil}{\lceil}{\rceil}
\DeclarePairedDelimiter\norm{\Vert}{\Vert}
\DeclarePairedDelimiter\enbrace{(}{)}
\DeclarePairedDelimiter\benbrace{[}{]}
\DeclarePairedDelimiter\bbenbrace{[\![}{]\!]}
\DeclarePairedDelimiter\lenbrace{<}{>}
\DeclarePairedDelimiter\angbrace{\langle}{\rangle}
\newcommand{\ssbrace}[1]{{\scriptscriptstyle\enbrace{#1}}}
\newcommand{\ssbbrace}[1]{{\scriptscriptstyle\benbrace{#1}}}
% ======================================================================================

%-- Mengen
% ======================================================================================
\newcommand\SetSymbol[1][]{\nonscript\:#1\vert\allowbreak\nonscript\:\mathopen{}}
\providecommand\given{} % to make it exist
\DeclarePairedDelimiterX\set[1]\{\}{\renewcommand\given{\SetSymbol[\delimsize]}#1}
% ======================================================================================

%-- Skalarprodukt (3 Varianten) 
% ======================================================================================
\DeclarePairedDelimiterX\sprod[2]{\langle}{\rangle}{#1\,\delimsize\vert\,#2}
\DeclarePairedDelimiterX\skal[2]{\langle}{\rangle}{#1\,,\,#2}
\makeatletter
\DeclareFontFamily{OMX}{MnSymbolE}{}
\DeclareSymbolFont{MnLargeSymbols}{OMX}{MnSymbolE}{m}{n}
\SetSymbolFont{MnLargeSymbols}{bold}{OMX}{MnSymbolE}{b}{n}
\DeclareFontShape{OMX}{MnSymbolE}{m}{n}{
    <-6>  MnSymbolE5
   <6-7>  MnSymbolE6
   <7-8>  MnSymbolE7
   <8-9>  MnSymbolE8
   <9-10> MnSymbolE9
  <10-12> MnSymbolE10
  <12->   MnSymbolE12
}{}
\DeclareFontShape{OMX}{MnSymbolE}{b}{n}{
    <-6>  MnSymbolE-Bold5
   <6-7>  MnSymbolE-Bold6
   <7-8>  MnSymbolE-Bold7
   <8-9>  MnSymbolE-Bold8
   <9-10> MnSymbolE-Bold9
  <10-12> MnSymbolE-Bold10
  <12->   MnSymbolE-Bold12
}{}
\let\llangle\@undefined
\let\rrangle\@undefined
\DeclareMathDelimiter{\llangle}{\mathopen}%
                     {MnLargeSymbols}{'164}{MnLargeSymbols}{'164}
\DeclareMathDelimiter{\rrangle}{\mathclose}%
                     {MnLargeSymbols}{'171}{MnLargeSymbols}{'171}
\makeatother
\DeclarePairedDelimiterX\sskal[2]{\llangle}{\rrangle}{#1\,,\,#2}
% ======================================================================================

%-- Abbildungsdefinition
% ======================================================================================
\newcommand{\mapdef}[5]{%
	\[
		\begin{array}{rcl}
			\textstyle #1 &\xrightarrow{\minwidthbox{#5}{2em}} & \textstyle #2 \\[0.5ex]
			\textstyle #3 &\xmapsto{\minwidthbox{\mbox{ }}{2em}} & \textstyle #4
		\end{array}
	\]
}
% ======================================================================================

%-- modifiziertes Stackrel 
% ======================================================================================
\newcommand{\StackText}[2]{\stackrel{\mbox{\scriptsize #1}}{#2}}
\newcommand{\StackTextClap}[2]{\stackrel{\mathclap{\mbox{\scriptsize #1}}}{#2}}
% ======================================================================================

%-- Blitz
% ======================================================================================
\newcommand{\light}{\text{\raisebox{-.3ex}{\Large\Lightning}}}
% ======================================================================================


%-- Underbrace u.Ä. als Befehl in LaTeX-Syntax (und ohne Spacingprobleme mit nachfolgenden Operatoren...)
% ======================================================================================
\newcommand{\Underbrace}[2]{{\underbrace{#1}_{#2}}}
\newcommand{\Underbracket}[2]{{\underbracket[0.7pt][2pt]{#1}_{#2}}}
\newcommand{\Overbracket}[2]{{\overbracket[0.7pt][2pt]{#1}^{#2}}}
% ======================================================================================


%-- Deklaration weiterer Operatoren (allgemein)
% ======================================================================================
\DeclareMathOperator{\re}{Re} % Realteil
\let\Re\relax
\DeclareMathOperator{\Re}{Re} % Realteil
\DeclareMathOperator{\im}{im} % Bild
\let\Im\relax
\DeclareMathOperator{\Im}{Im} % Bild
\DeclareMathOperator{\id}{id} % identische Abbildung
\DeclareMathOperator{\conj}{conj} % Konjugation
\DeclareMathOperator{\sgn}{sgn} % Signum
\DeclareMathOperator{\End}{End} % Endomorphismen
\DeclareMathOperator{\Hom}{Hom} % Homomorphismen
\DeclareMathOperator{\Iso}{Iso} % Isomorphismen
\DeclareMathOperator{\Aut}{Aut} % Automorphismen
\DeclareMathOperator{\Span}{span} % Span
\DeclareMathOperator{\coker}{coker} % Kokern
\DeclareMathOperator{\Tr}{Tr} % Spur,Trace
\DeclareMathOperator{\pr}{pr} % Projektion
\DeclareMathOperator{\diag}{diag} % Diagonalmatrix
\DeclareMathOperator{\Rg}{Rg} % Rang
\DeclareMathOperator{\const}{const} % konstante Abbildung
\DeclareMathOperator{\Spur}{Spur} % Spur
\DeclareMathOperator{\Arg}{Arg} % Argument
\DeclareMathOperator{\dist}{dist} % Distanz
\DeclareMathOperator{\supp}{supp} % Träger
\DeclareMathOperator{\Char}{char} % Charakteristik
% ======================================================================================


%-- Deklaration weiterer Operatoren (Differentiale etc.)
% ======================================================================================
\DeclareMathOperator{\grad}{grad} % Gradient
\DeclareMathOperator{\dive}{div} % Gradient
\DeclareMathOperator{\rot}{rot} % Rotation
\newcommand{\D}{\ensuremath{\mathrm{D}\mkern-1.0mu}} % Differential
\newcommand{\mathd}{\ensuremath{\mathrm{d}\mkern-1.0mu}} % äußere Ableitung
\newcommand{\Tmap}{\ensuremath{\mathrm{T}\mkern-0.85mu}} % Tangentialraum
\let\Tang\Tmap
\DeclareMathOperator{\Diff}{Diff}
\newcommand{\diff}[2]{\ensuremath{\frac{{\partial #1}}{{\partial #2}} }}
\newcommand{\diffd}[2]{\ensuremath{\frac{\mathd #1}{\mathd #2} }}
% ======================================================================================


%-- Deklaration weiterer Operatoren (Topologie)
% ======================================================================================
\newcommand*\interior[1]{\overset{\smash{\raisebox{-0.18ex}{$\scriptstyle\circ$}}}{#1}}
\newcommand{\sing}{{\raisemath{1.1pt}{\scriptscriptstyle\mathrm{sing}}}}
\newcommand{\pt}{\mathrm{pt}}
\DeclareMathOperator{\Zyl}{Zyl}
\DeclareMathOperator{\Tel}{Tel}
\newcommand{\op}{\mathrm{op}}
\DeclareMathOperator{\Sp}{Sp}
\DeclareMathOperator{\Keg}{Keg}
\newcommand{\slashedi}{i\hspace{-3.5pt}/}
\newcommand{\cupp}{\smallsmile}
\newcommand{\capp}{\smallfrown}
\DeclareMathOperator*{\colim}{colim}
\DeclareMathOperator{\PD}{PD}
\newcommand{\lf}{\mathrm{lf}}
\DeclareMathOperator{\sig}{sig}
\DeclareMathOperator{\Tor}{Tor}
\DeclareMathOperator{\Ext}{Ext}
\DeclareMathOperator{\AW}{AW}
\DeclareMathOperator{\Proj}{Proj}
\DeclareMathOperator{\Gr}{Gr}
\DeclareMathOperator{\res}{res}
\DeclareMathOperator{\Spec}{Spec}
\DeclareMathOperator{\co}{co}
\DeclareMathOperator{\ch}{ch}
\DeclareMathOperator{\wOp}{w}
\DeclareMathOperator{\Ar}{Ar}
\newcommand{\actson}{\mathrel{\curvearrowright}}
\let\acts\actson
\let\action\actson
\DeclareMathSymbol{\bbDelta}{\mathord}{bbold}{"01}
\newcommand{\DDelta}{\bbDelta}
\DeclareMathOperator{\Star}{Star}
\DeclareMathOperator{\Link}{Link}
\DeclareMathOperator{\EPK}{EPK}
\DeclareMathOperator{\Vol}{Vol}
\newcommand{\cell}{{\raisemath{1.1pt}{\scriptscriptstyle\mathrm{cell}}}}
\DeclarePairedDelimiter{\homologieklasse}{\llbracket}{\rrbracket}
\newcommand{\rand}[1]{\ensuremath{\partial^{\scriptscriptstyle #1}}}
\DeclareMathOperator{\ab}{ab}
% ======================================================================================


%-- Deklaration von Operatoren (Liegruppen)
% ======================================================================================
\DeclareMathOperator{\GL}{GL}
\DeclareMathOperator{\SO}{SO}
\DeclareMathOperator{\Ad}{Ad}
\DeclareMathOperator{\ad}{ad}
\DeclareMathOperator{\On}{O}
\DeclareMathOperator{\Un}{U}
\DeclareMathOperator{\SU}{SU}
\DeclareMathOperator{\Mat}{Mat}
\DeclareRobustCommand{\Der}{\mathop{\mathfrak{der}}}
\DeclareMathOperator{\SL}{SL}
\DeclareMathOperator{\Graph}{Graph}
\DeclareMathOperator{\Int}{Int}
\DeclareRobustCommand{\intAlg}{\mathop{\mathfrak{int}}}
\DeclareMathOperator{\aut}{aut}
\DeclareMathOperator{\Rad}{Rad}
\DeclareMathOperator{\Nil}{Nil}
\DeclareMathOperator{\rad}{rad}
\DeclareMathOperator{\nil}{nil}
\DeclareMathOperator{\Ric}{Ric}
\DeclareMathOperator{\ric}{ric}
\newcommand{\bi}{\mathrm{bi}}
\DeclareMathOperator{\Isom}{Isom}
\DeclareMathOperator{\Sym}{Sym}
\newcommand{\opL}{\ensuremath{\mathrm{L}\mkern-0.6mu}}
% ======================================================================================

%-- Deklaration von Operatoren (Funktionalanalysis)
% ======================================================================================
\DeclareMathOperator{\tr}{tr}
\newcommand{\w}{\mkern1mu\mathrm{w}}
\newcommand{\sa}{\mathrm{sa}}
\newcommand{\vb}{\mathrm{v\mkern-2.5mu.b\mkern-1.5mu.}} % vollständig beschränkt
\newcommand{\so}{\mathrm{\mkern.3mu s\mkern-1.4mu.\mkern-.6mu o\mkern-1.7mu.}} % \newcommand{\so}{\mathrm{s.o.}}
\newcommand{\solim}{\so\text{-}\mkern-0.8mu\lim}
\newcommand{\wo}{\mathrm{w\mkern-3mu.\mkern-.4mu o\mkern-1.7mu.}}
\newcommand{\Top}[1]{\mathcal{T}_{\mkern-2.3mu #1}}
\newcommand{\weakT}[1]{\ensuremath{\mathcal{T}_{#1}^{\mkern+1.0mu\text{\raisebox{0.4ex}{$\mathrm{w}$}}}}}
\newcommand{\weakTstar}[1]{\ensuremath{\mathcal{T}_{#1}^{\mkern+1.0mu\text{\raisebox{0.4ex}{$\mathrm{w}$}}^*}}}
\newcommand{\TWeakStar}{\Top{\w^*}}
\newcommand{\TWeakOp}{\Top{\wo}}
\newcommand{\Tso}{\Top{\so}}
\newcommand{\finSub}{\subset\mkern-0.7mu \subset}
\DeclareMathOperator{\Inv}{Inv}
\newcommand{\simm}{{\hspace{-1.6pt}\raisemath{0.5pt}{\sim}}}
\newcommand{\plus}{{\hspace{-1.6pt}+}}
\DeclareMathOperator{\ev}{ev}
\DeclareMathOperator{\Alg}{Alg}
\DeclareMathOperator{\her}{her}
\newcommand{\subher}{\subset_{\her}}
\newcommand{\grenzw}[1]{\xrightarrow{\minwidthbox{#1}{1.4em}}}
\newcommand{\grenzwl}[1]{\xleftarrow{\minwidthbox{#1}{1.4em}}}
\newcommand{\grenzwIn}[1]{\grenzw{\raisemath{-2pt}{#1}}}
\newcommand{\MyTo}[1]{\tikzexternaldisable\mathbin{\tikz[baseline] \draw[-to,line width=.4pt] (0ex,0.94ex) -- (#1,0.94ex);}\tikzexternalenable}
\newcommand{\dlim}{%
    \mathchoice
      {\lim\limits_{\MyTo{4.2ex}}}% \displaystyle
      {\lim\limits_{\MyTo{2.8ex}}}% \textstyle
      {\lim\limits_{\MyTo{2.3ex}}}% \scriptstyle
      {\lim\limits_{\MyTo{2.3ex}}}% \scriptscriptstyle
}
\newcommand{\Dlim}{\killDescendersM{\dlim}}
\DeclareMathOperator{\sep}{sep}
\DeclareMathOperator{\diam}{diam}
\DeclareMathOperator{\conv}{conv}
\DeclareMathOperator{\Prim}{Prim}
\DeclareMathOperator{\hull}{hull}
\DeclareMathOperator{\red}{red}
\DeclarePairedDelimiterX\bra[1]{\langle}{\rvert}{#1\,}
\DeclarePairedDelimiterX\ket[1]{\lvert}{\rangle}{\,#1}
\DeclarePairedDelimiterX\bracket[2]{\langle}{\rangle}{#1\,\delimsize\vert\,#2}
\newcommand{\tensormax}{\mathbin{\otimes_{\max}}}
\newcommand{\tensormin}{\mathbin{\otimes_{\min}}}
\DeclareMathOperator{\Ped}{Ped}
\newcommand{\alg}{\mathrm{alg}}
\DeclareMathOperator{\CPC}{CPC}
\DeclareMathOperator{\CP}{CP}
\DeclareMathOperator{\UPC}{UPC}
\newcommand{\DeltaOp}{\mathbin{\Delta}}
\newcommand{\kernedP}{\mathcal{P}\mkern-2mu}
\newcommand{\Pinfty}{\kernedP_{\infty}}
\DeclareMathOperator{\Groth}{Groth}
\DeclareMathOperator{\rk}{rk}
\newcommand{\MvN}{\mathrm{MvN}}
% ======================================================================================

%-- Kategorien
% ======================================================================================
\DeclareMathOperator{\Mor}{Mor}
\DeclareMathOperator{\Obj}{Obj}
\newcommand{\TOP}{\textsc{Top}}
\newcommand{\HTOP}{\textsc{HTop}}
\newcommand{\VR}{\textsc{VR}}
\newcommand{\MOD}{\textsc{Mod}}
\newcommand{\Mod}[1]{#1\text{-}\MOD}
\newcommand{\MONOIDE}{\textsc{Monoide}}
\newcommand{\SET}{\textsc{Set}}
\newcommand{\MAN}{\textsc{Man}}
\newcommand{\GRUPPEN}{\textsc{Gruppen}}
\newcommand{\ABELGRUPPEN}{\textsc{Abel.Gruppen}}
\newcommand{\ABEL}{\textsc{Abel}}
\newcommand{\KAT}{\textsc{Kat}}
\newcommand{\FUN}{\textsc{Fun}}
\newcommand{\SIMP}{\textsc{Simp}}
\newcommand{\VEKT}{\textsc{Vekt}}
\newcommand{\CH}{\textsc{Ch}}
\newcommand{\CSTARUN}{C^*\text{-}\textsc{Alg}^{\raisemath{-2.5pt}{1}}}
\newcommand{\CSTAR}{C^*\text{-}\textsc{Alg}}
\newcommand{\AB}{\textsc{Ab}}
% ======================================================================================
% ======================================================================================



% -- theorem packages
% ======================================================================================
\usepackage{amsthm}
\usepackage{thmtools,thm-restate}
\usepackage{mdframed}
\renewcommand{\listtheoremname}{Übersicht aller Aussagen}
\usepackage{bookmark}
\bookmarksetup{open,numbered}
\makeatletter
\newcommand*{\theorembookmark}{%
  \bookmark[
    dest=\@currentHref,
    rellevel=1,
    keeplevel,
  ]{%
    \thmt@thmname\space\csname the\thmt@envname\endcsname
    \ifx\thmt@shortoptarg\@empty
    \else
      \space(\thmt@shortoptarg)%
    \fi
  }%
}   
\makeatother
% ======================================================================================

% -- Definition der einzelnen Theorem-Umgebungen
% ======================================================================================
\declaretheoremstyle[%
	headfont=\sffamily\bfseries,
	notefont=\normalfont\sffamily\scshape,
	bodyfont=\normalfont,
	headformat=\NUMBER\ \NAME\NOTE,
	headpunct=.,
	postheadspace=1em,
	spaceabove=15pt,spacebelow=10pt,
	shaded={bgcolor=gray!20},
	postheadhook=\theorembookmark]%
{mainstyle}
\declaretheoremstyle[%
	headfont=\sffamily\bfseries,
	notefont=\normalfont\sffamily\scshape,
	bodyfont=\normalfont,
	headformat=\NUMBER\ \NAME\NOTE,
	headpunct=.,
	postheadspace=1em,
	spaceabove=15pt,spacebelow=10pt,
	shaded={bgcolor=fb10_blue!20},
	postheadhook=\theorembookmark]%
{mainstyle_blue}
\declaretheoremstyle[%
	headfont=\sffamily\bfseries,
	notefont=\normalfont\sffamily\scshape,
	bodyfont=\normalfont,
	headformat=\NUMBER\ \NAME\NOTE,
	headpunct=.,
	postheadspace=1em,
	spaceabove=15pt,spacebelow=10pt,
	postheadhook=\theorembookmark]%
{mainstyle_unshaded}
\declaretheoremstyle[%
	headfont=\sffamily\bfseries,
	notefont=\normalfont\sffamily\scshape,
	bodyfont=\normalfont,
	headformat=\NUMBER\NAME\NOTE,
	headpunct=.,
	postheadspace=1em,
	spaceabove=15pt,spacebelow=10pt,
	% shaded={bgcolor=gray!20},
	postheadhook=\theorembookmark]%
{mainstyle_unnumbered}
\declaretheoremstyle[%
	headfont=\sffamily\bfseries,
	notefont=\normalfont\sffamily\scshape,
	bodyfont=\normalfont,
	headformat=swapnumber,
	headpunct=.,
	postheadspace=1em,
	spaceabove=15pt,spacebelow=10pt,
	shaded={bgcolor=gray!20},
	postheadhook=\theorembookmark,
	qed=\qedsymbol]%
{mainstyleB}
\declaretheoremstyle[%
	headfont=\bfseries\scshape,
	bodyfont=\normalfont,
	headpunct=:,
	postheadspace=1em,
	spacebelow=12pt,spaceabove=2pt,
	qed=\qedsymbol]%
{beweise}
\declaretheoremstyle[%
	headfont=\bfseries\scshape,
	bodyfont=\normalfont,
	headpunct=:,
	postheadspace=1em,
	spacebelow=12pt,spaceabove=2pt]%
{beweisskizze}
\declaretheoremstyle[%
	headfont=\sffamily\bfseries,
	bodyfont=\normalfont,
	headpunct=:,
	postheadspace=1em,
	spacebelow=10pt,spaceabove=10pt]%
{bemerkungen}
\declaretheorem[name=Definition,parent=section,style=mainstyle_blue]{definition}
\declaretheorem[name=Definition \& Proposition,refname=Proposition,sharenumber=definition,style=mainstyle_blue]{definitionP}
\declaretheorem[name=Definition,numbered=no,style=mainstyle_unnumbered]{definition*}
\declaretheorem[name=Theorem,sharenumber=definition,style=mainstyle]{theorem}
\declaretheorem[name=Theorem,numbered=no,style=mainstyle_unnumbered]{theorem*}
\declaretheorem[name=Proposition,sharenumber=definition,style=mainstyle,refname=Proposition]{proposition}
\declaretheorem[name=Lemma,sharenumber=definition,style=mainstyle]{lemma}
\declaretheorem[name=Satz,sharenumber=definition,style=mainstyle,refname=Satz]{satz}
\declaretheorem[name=Satz,sharenumber=definition,style=mainstyle_unshaded]{satzUnshaded}
\declaretheorem[name=Definition,sharenumber=definition,style=mainstyle_unshaded]{definitionUnshaded}
\declaretheorem[name=Satz,numbered=no,style=mainstyle_unnumbered]{satz*}
\declaretheorem[name=Korollar,sharenumber=definition,style=mainstyle,refname=Korollar]{korollar}
\declaretheorem[name=Korollar,sharenumber=definition,style=mainstyleB,refname=Korollar]{korollarB}
\declaretheorem[name=Frage,numbered=no,style=mainstyle_unnumbered]{frage}
\declaretheorem[name=Frage,sharenumber=definition,style=mainstyle_unshaded]{frageA}
\declaretheorem[name=Erinnerung,sharenumber=definition,style=mainstyle_unshaded]{erinnerungA}
\declaretheorem[name=Ausblick,sharenumber=definition,style=mainstyle_unshaded]{ausblick}
\declaretheorem[name=Konvention,sharenumber=definition,style=mainstyle]{konvention}
\declaretheorem[name=Notation,sharenumber=definition,style=mainstyle_unshaded]{notation}
\declaretheorem[name=Bemerkung,sharenumber=definition,style=mainstyle_unshaded,refname=Bemerkung]{bemerkung}
\declaretheorem[name=Bemerkung,numbered=no,style=mainstyle_unnumbered]{bemerkung*}
\declaretheorem[name=Beispiel,sharenumber=definition,style=mainstyle_unshaded,refname=Beispiel]{beispiel}
\declaretheorem[name=Beispiel,numbered=no,style=mainstyle_unnumbered]{beispiel*}
\declaretheorem[name=Exkurs,numbered=no,style=mainstyle_unnumbered]{exkurs*}
\declaretheorem[name=Beweis,numbered=no,style=beweise]{beweis}
\declaretheorem[name=Übung,numbered=no,style=bemerkungen]{uebung}
\declaretheorem[name=Erinnerung,numbered=no,style=bemerkungen]{erinnerung}

% english versions
\declaretheorem[name=Remark,sharenumber=definition,style=mainstyle_unshaded]{remark}
\declaretheorem[name=Remark,numbered=no,style=mainstyle_unnumbered]{remark*}
\declaretheorem[name=Example,sharenumber=definition,style=mainstyle_unshaded]{example}
\declaretheorem[name=Corollary,sharenumber=definition,style=mainstyle]{corollary}
\let\proof\relax
\declaretheorem[name=Proof,numbered=no,style=beweise]{proof}
\declaretheorem[name=Sketch of Proof,numbered=no,style=beweisskizze]{sketch}
% ======================================================================================

%--Inhaltsverzeichnis
% ======================================================================================
\usepackage[tocindentauto]{tocstyle}
\usetocstyle{KOMAlike}
% ======================================================================================

%-- Dinge, die erst am Ende getan werden dürfen
% ======================================================================================
\shorthandon{"}
\usepackage{ellipsis}
% ======================================================================================


\newcommand{\fach}{Operatoralgebren \RM{2}.}
\newcommand{\semester}{Sose 2016}
\newcommand{\homepage}{http://wwwmath.uni-muenster.de/u/wilhelm.winter/wwinter/operatoralgebren_II.html}

\newcommand{\prof}{Prof.\ Dr.\ Wilhelm Winter}
\publishers{\scalebox{10}{\Huge$^{\raisebox{2.2pt}{\text{\small\RM{2}.}}}\mkern-1.5mu C^*$}}
\providecommand{\verfasser}{Jannes Bantje}

%--Konfiguration von scrheadings
\setheadsepline{1pt}[\color{light_gray}]
\pagestyle{scrheadings}
%\fancyhf[H, F]{}
\clearscrheadfoot

\providecommand{\shortFach}{\fach}
\lehead{\includegraphics[height=0.6 cm,keepaspectratio]{../!config/Bilder/Logo_WWU_Muenster_light_gray.pdf}}
\rehead{\rule{0cm}{0.6cm}\footnotesize\sffamily\color{light_gray}\verfasser{} -- Mitschrift \shortFach}
\lohead{\rule{0cm}{0.6cm}\footnotesize\sffamily\color{light_gray}Stand: \today \; \thistime[:]}
\rohead{\includegraphics[height=0.6 cm,keepaspectratio]{../!config/Bilder/fb10logo_gray.pdf}}


\ofoot[{ \color{dark_gray} \LARGE \sffamily \thepage}]{{ \color{dark_gray} \LARGE \sffamily \thepage}} %hier wir auch der plain Stil bearbeitet!
\automark{section}
\ifoot{ \color{dark_gray} \small \leftmark}

%--Metadaten
\providecommand{\mail}{j.bantje@wwu.de}
\author{\verfasser}
\titlehead{\includegraphics[height=1.5cm, keepaspectratio]{../!config/Bilder/Logo_WWU_Muenster.pdf}%
\hfill \includegraphics[height=1.3cm, keepaspectratio]{../!config/Bilder/fb10logo.pdf}}
\title{Skript \fach}
\subtitle{Mitschrift der Vorlesung  \enquote{\fach} von \prof}

\KOMAoptions{DIV=last}


\begin{document}
\pagenumbering{Roman}
\maketitle
\begin{abstract}
\section*{Aktuelle Version verfügbar bei}
\newcommand{\dieBreite}{11cm}
\begin{minipage}{4cm}
	\qrcode[height=3.3cm, version=6]{https://gitlab.com/JaMeZ-B/LaTeX-WWU}
\end{minipage}
\hfill
\begin{minipage}{\dieBreite}
	% \includegraphics[height=0.6cm, keepaspectratio]{../!config/Bilder/wm_no_bg.pdf}
	\includegraphics[height=0.8cm, keepaspectratio]{../!config/Bilder/wm_no_bg.pdf}\\
	\url{https://gitlab.com/JaMeZ-B/LaTeX-WWU} \smallskip\\
	Das zentrale Repository des \enquote{\LaTeX-WWU}-Projekts befindet sich auf der Plattform GitLab.com.
	Neben der Koordination aller Beteiligten werden über diesen Dienst auch die PDFs gebaut, die in der Readme verlinkt sind.
\end{minipage}\\[1cm]
\begin{minipage}{4cm}
	\qrcode[height=3.3cm, version=6]{https://github.com/JaMeZ-B/latex-wwu}
\end{minipage}
\hfill
\begin{minipage}{\dieBreite}
	\includegraphics[height=0.6cm, keepaspectratio]{../!config/Bilder/github_octo.pdf}
	\includegraphics[height=0.6cm, keepaspectratio]{../!config/Bilder/GitHub_Logo.pdf}\\
	\url{https://github.com/JaMeZ-B/latex-wwu} \smallskip\\
	Die Entwicklung des \enquote{\LaTeX-WWU}-Projekts hat ursprünglich auf GitHub stattgefunden, ist mittlerweile aber zu GitLab gewechselt.
	Das GitHub-Repository wird stündlich automatisch aktualisiert, Merge-Requests werden aber nicht mehr entgegengenommen.
\end{minipage}\\[1cm]
% \begin{minipage}{4cm}
% 	\qrcode[height=3.3cm, version=6]{https://uni-muenster.sciebo.de/public.php?service=files&t=965ae79080a473eb5b6d927d7d8b0462}
% \end{minipage}
% \hfill
% \begin{minipage}{\dieBreite}
% 	\raisebox{-2pt}{\includegraphics[height=0.6cm, keepaspectratio]{../!config/Bilder/sciebo_logo.pdf}}
% 	\resizebox{!}{0.5cm}{\large \sffamily\textbf{sciebo}} {\sffamily\large die Campuscloud} \\
% 	\resizebox{\dieBreite}{!}{\footnotesize\url{https://uni-muenster.sciebo.de/public.php?service=files&t=965ae79080a473eb5b6d927d7d8b0462}}\smallskip\\
% 	Sciebo ist ein Dropbox-Ersatz der Hochschulen in NRW, der von der Uni Münster in leitender Position auf Basis der OpenSource-Software Owncloud aufgebaut wurde.
% \end{minipage}\\[1cm]
\hrule \mbox{ }\\[0.7cm]
\begin{minipage}{4cm}
	\qrcode[height=3.3cm, version=6]{\homepage}
\end{minipage}
\hfill
\begin{minipage}{\dieBreite}
	\resizebox{!}{0.5cm}{\large\sffamily\textbf{Vorlesungshomepage}}\\
	\resizebox{\dieBreite}{!}{\footnotesize\url{\homepage}}\smallskip\\
	Hier ist ein Link zur offiziellen Vorlesungshomepage.
\end{minipage}
\newpage
\section*{Vorwort --- Mitarbeit am Skript}
Dieses Dokument ist eine Mitschrift aus der Vorlesung \enquote{\fach, \semester}, gelesen von \prof. 
Der Inhalt entspricht weitestgehend dem Tafelanschrieb. 
Für die Korrektheit des Inhalts übernehme ich keinerlei Garantie! 
Für Bemerkungen und Korrekturen -- und seien es nur Rechtschreibfehler -- bin ich sehr dankbar. 
Korrekturen lassen sich prinzipiell auf drei Wegen einreichen: 
\begin{itemize}
	\item Persönliches Ansprechen in der Uni, Mails an \hrefsymmail{mailto:\mail}{\mail} (gerne auch mit annotieren PDFs) oder Kommentare auf \url{https://gitlab.com/JaMeZ-B/LaTeX-WWU}.
	\item \emph{Direktes} Mitarbeiten am Skript: Den Quellcode poste ich auf GitLab (siehe oben), also stehen vielfältige Möglichkeiten der Zusammenarbeit zur Verfügung:
	Zum Beispiel durch Kommentare am Code über die Website und die Kombination Fork und Merge-Request. 
	Wer sich verdient macht oder ein Skript zu einer Vorlesung, die ich nicht besuche, beisteuern will, dem gewähre ich gerne auch Schreibzugriff.
	
	Beachten sollte man dabei, dass dazu ein Account bei \url{gitlab.com} notwendig ist, der allerdings ohne Angabe von persönlichen Daten angelegt werden kann. 
	Wer bei GitLab (bzw. dem zugrunde liegenden Open-Source-Programm \enquote{\texttt{git}}) -- verständlicherweise -- Hilfe beim Einstieg braucht, dem helfe ich gerne weiter. 
	Es gibt aber auch zahlreiche empfehlenswerte Tutorials im Internet.\footnote{zB. \url{https://try.github.io/levels/1/challenges/1}, ist auf Englisch, aber dafür interaktiv}
	\item \emph{Indirektes} Mitarbeiten: \TeX-Dateien per Mail verschicken. 
	
	Dies ist nur dann sinnvoll, wenn man einen ganzen Abschnitt ändern möchte (zB. einen alternativen Beweis geben), da ich die Änderungen dann per Hand einbauen muss! Ich freue mich aber auch über solche Beiträge!
\end{itemize}
\section*{Anmerkung}
Innerhalb dieser Mitschrift wird man öfter den Ausdruck \enquote{Warum?} finden. Dies sind vom Dozenten bewusst weggelassene Details, die zu verstärktem Mitdenken beim Lesen 
animieren sollen. Oftmals sind dies schon aus vorherigen Semestern bekannte Sachverhalte. Nur an wenigen Stellen habe ich die fehlenden Details hinzugefügt.

\section*{Literatur}
\begin{itemize}[itemsep=0pt]
	\item G. J. \textsc{Murphy} \emph{C*-algebras and Operator Theory} \cite{Murphy}
	\item N. P. \textsc{Brown}, N. \textsc{Ozawa} \emph{C*-algebras and Finite-Dimensional Approximations} \cite{brown2008}
	\item M. \textsc{Rørdam}, F. \textsc{Larsen}, N. J. \textsc{Laustsen} \emph{An Introduction to K-theory for C*-Algebras} \cite{rordam2000}
	\item N. E. \textsc{Wegge-Olsen} \emph{K-theory and C*-algebras} \cite{wegge}
	\item B. \textsc{Blackadar} \emph{K-theory for Operator Algebras} \cite{blackadar1998}
	\item B. \textsc{Blackadar} \emph{Operator Algebras} \cite{blackadar2006}
	\item M. \textsc{Takesaki} \emph{Theory of Operator Algebras I.} \cite{takesaki2002}
	\item R.V. \textsc{Kadison}, J. R. \textsc{Ringrose} \emph{Fundamentals of the Theory of Operator Algebras Volumes I. \& II.} \cite{kadison1997-1} \& \cite{kadison1997-2}
	\item G.K. \textsc{Pedersen} \emph{C*-algebras and their Automorphism Groups} \cite{PedC*}
\end{itemize}
\end{abstract}

\tableofcontents
\cleardoubleoddemptypage

\pagenumbering{arabic}
\setcounter{page}{1}
\setcounter{footnote}{0}

\section{Tensorprodukte} % (fold)
\label{sec:1}

\begin{erinnerungA}[{name=[Tensorprodukt]}]
	Seien $V,W$ $\mathbb{K}$-Vektorräume.
	Sei $X$ ein weiterer $\mathbb{K}$-Vektorraum und $\phi \colon V\times W \to X$ bilinear.
	Dann heißt $(X,\varphi)$ \Index{Tensorprodukt} von $V$ und $W$, falls folgende universelle Eigenschaft gilt:
	\[
		\begin{tikzcd}[sep=large]
			V \times W \drar["\text{bil.}"'] \rar["\varphi"] & X \dar[dashed,"\exists ! \psi \text{ lin.} "] \\
			& U
		\end{tikzcd}
	\]
\end{erinnerungA}

\begin{definitionP}[label=propdef:12,{name=[{Konstruktion des algebraischen Tensorproduktes}]}]
	$(X,\varphi)$ wie oben existiert und ist eindeutig bis auf linearen Isomorphismus.
	Wir schreiben $V \odot W$ für $X$ und $v \otimes w$ für $\varphi(v,w)$.\marginnote{$\odot$ algebraisches Tensorprodukt}
\end{definitionP}
\begin{beweis}
	Sei $M$ der freie Vektorraum mit Basis $V \times W$, anders ausgedrückt
	\[
		M:= \set[\big]{f \colon V\times W \to \mathbb{K} \given \supp f \text{ endlich}}
	\]
	Sei nun 
	\[
		N := \Span \set*{
			\begin{array}{l}
				(v_1 +v_2,w) - \enbrace[\big]{(v_1,w) + (v_2,w)} \\
				(v,w_1+w_2) - \enbrace[\big]{(v,w_1) + (v,w_2)} \\
				(v,\lambda \cdot w) - \lambda \cdot (v,w), (\lambda \cdot v,w) - \lambda \cdot (v,w) 
			\end{array}
		\given 
		\begin{array}{l}
			v,v_1,v_2 \in V,\\
			w,w_1,w_2 \in W,\\
			\lambda \in \mathbb{K}
		\end{array}}
	\]
	Wir setzen $X := \sfrac{M}{N}$ und $\varphi(v,w) := (v,w)+N$.
	Der Rest ist eine einfache Übung.
\end{beweis}

\begin{proposition}[{name=[Tensorprodukt von Algebren]}]
	Seien $A,B$ $\mathbb{K}$-Algebren.
	Dann ist $A \odot B$ wieder eine $\mathbb{K}$-Algebra mit 
	\[
		(a \otimes b) (a' \otimes b') = aa' \otimes bb'
	\]
	Sind $A$ und $B$ $^*$-Algebren, dann ist $A \odot B$ wieder eine $^*$-Algebra mit 
	\[
		(a \otimes b)^* = a^* \otimes b^*
	\]
	Sind $A$ und $B$ kommutativ/unital, so auch $A \odot B$.
\end{proposition}
\begin{beweis}
	Betrachte $M$ wie in \autoref{propdef:12} und überprüfe, dass Multiplikation auf $\sfrac{M}{N}$ wohldefiniert und assoziativ ist.
	
	Alternativ: Die Elemente $(a,b) \in A \times B$ liefern eine bilineare Abbildung
	\mapdef{M_{(a,b)} \colon A \times B}{A \odot B}{(a',b')}{aa' \otimes bb'}{}
	Aus der universellen Eigenschaft erhält man $m_{(a,b)} \colon A \odot B \to A \odot B$ linear.
	Dies liefert wiederum eine bilineare Abbildung $m \colon A \times B \to \mathcal{L}(A \odot B, A \odot B)$ durch $(a,b) \mapsto m_{(a,b)}$.
	Nach der universellen Eigenschaft bekommt man nun eine lineare Abbildung $A \odot B \to \mathcal{L}(A \odot B, A\odot B)$.
	Daraus erhält man eine bilineare Abbildung
	\[
		\mu \colon A \odot B \times A \odot B \longrightarrow A \odot B
	\]
	Dann ist $(A \odot B,\mu)$ eine $\mathbb{K}$-Algebra.
\end{beweis}

\begin{bemerkung}[label=bem:14,{name=[{universelle Eigenschaft des Tensorproduktes von *-Algebren}]}]
	\begin{enumerate}[(i)]
		\item Seien $A,B$ $^*$-Algebren.
		Dann hat das Tensorprodukt $A \odot B$ als $^*$-Algebra die folgende universelle Eigenschaft:
		
		Falls $C$ eine $^*$-Algebra ist und $\pi_A \colon A \to C$, $\pi_B \colon B \to C$ $^*$-Homomorphismen sind, sodass für den Kommutator $\benbrace*{\pi_A(A), \pi_B(B)}=0$ gilt,\marginnote{$[A,B]$ besteht aus den Elementen der Form $ab -ba$} so existiert genau ein $^*$-Homomorphismus $\pi \colon A \odot B \to C$ mit $\pi(a \otimes b)= \pi_A(a) \cdot \pi_B(b)$ für alle $a \in A$, $b \in B$.
		\item Es seien $A,B,C,D$ $^*$-Algebren, $\pi_A \colon A \to C$ und $\pi_B \colon B \to D$ $^*$-Homomorphismen.
		Dann existiert genau ein $\pi \colon A \odot B \to C \odot D$ mit $\pi(a \otimes b)= \pi_A(a) \otimes \pi_B(b)$.
	\end{enumerate}
\end{bemerkung}

\begin{definitionP}[{name=[{Tensorprodukt von Hilberträumen}]}]
	Es seien $\mathcal{H}_1$, $\mathcal{H}_2$ Hilberträume.
	Dann ist $\mathcal{H}_1 \odot \mathcal{H}_2$ ein Prä-Hilbertraum mit
	\[
		\skal*{\xi \otimes \eta}{\xi' \otimes \eta'} := \skal*{\xi}{\xi'} \cdot \skal*{\eta}{\eta'} 
	\]
	Das \bet{Tensorprodukt von Hilberträumen}\index{Tensorprodukt!von Hilberträumen} ist nun definiert durch
	\[
		\mathcal{H}_1 \otimes \mathcal{H}_2 := \overline{\mathcal{H}_1 \odot \mathcal{H}_2}^{\skal*{\cdot}{\cdot}}
	\]
\end{definitionP}
\begin{beweis}
	Es seien $\xi \in \mathcal{H}_1$ und $\eta \in \mathcal{H}_2$.
	Nach der universellen Eigenschaft existiert eine lineare Abbildung $\tau_{(\xi,\eta)} \colon \mathcal{H}_1 \odot \mathcal{H}_2 \to \mathbb{C}$ mit $\tau_{{\xi,\eta}} \enbrace*{\xi' \otimes \eta'} = \skal*{\xi}{\xi'} \cdot \skal*{\eta}{\eta'}$.
	Die Abbildung $(\xi,\eta) \mapsto \overline{\tau_{{\xi,\eta}}}$ ist dann bilinear.
	Aus der universellen Eigenschaft erhalten wir eine antilineare Abbildung $\xi \otimes \eta \mapsto \tau_{(\xi,\eta)}$.
	Daraus erhält man eine Sesquilinearform 
	\[
		\skal*{\cdot}{\cdot } \colon \mathcal{H}_1 \odot \mathcal{H}_2 \longrightarrow \mathbb{C}
	\]
	Wir müssen noch zeigen, dass diese positiv definit ist:
	Sei $z = \sum_{i=1}^{n} \xi_i \otimes \eta_i \in \mathcal{H}_1 \odot \mathcal{H}_2$.
	Wähle eine Orthonormalbasis $\zeta_1, \ldots , \zeta_m$ für $\Span \set*{\eta_1, \ldots,\eta_n}$.
	Dann ist $z=\sum_{j=1}^{m} \xi_j' \otimes \zeta_j$ für $\xi_j' \in \mathcal{H}_1$ geeignet.
	Damit folgt 
	\[
		\skal*{z}{z} = \sum_{j=1}^{m} \norm*{\xi'_j}^2 \ge 0
	\]
	Insbesondere ist $z=0$ genau dann, wenn $\xi_1' = \ldots = \xi_m'=0$ ist, also $\skal*{z}{z}=0$ gilt.
\end{beweis}

\begin{bemerkung}[{name=[{Einbettung der Tensorprodukte beschränkter Operatoren}]},label=bem:16]
	Es existiert eine Einbettung $\iota \colon \mathcal{B}(\mathcal{H}_1) \odot \mathcal{B}(\mathcal{H}_2) \hookrightarrow \mathcal{B}(\mathcal{H}_1 \otimes \mathcal{H}_2)$ gegeben durch $(a \otimes b)(\xi \otimes \eta) = a(\xi) \otimes b (\eta)$. \emph{Siehe \cref{sub:injek_tensor_beschraenkte_operatoren}}
\end{bemerkung}

\begin{frageA}[{name=[{C*-Normen auf dem algebraischen Tensorprodukt}]}]
	Seien $A$ und $B$ $C^*$-Algebren.
	Die spannende Frage ist nun:
	\begin{center}
		\fbox{{\large¿}\,Was für $C^*$-Normen existieren auf $A \odot B$\,{\large?}}
	\end{center}
\end{frageA}

\begin{bemerkung}[{name=[Tensorprodukt aus C*-Norm auf dem algebraischen Tensorprodukt]}]
	Falls $\gamma \colon A \odot B \to \mathbb{R}^+$ eine $C^*$-Norm ist, so ist $A \otimes_\gamma B := \overline{A \odot B}^\gamma$ eine $C^*$-Algebra.\todo{RevChap 1}
\end{bemerkung}

\begin{proposition}[{name=[{minimale C*-Norm auf dem Tensorprodukt}]},label=prop:19]
	Es seien $A$ und $B$ $C^*$-Algebren.
	\begin{enumerate}[(i)]
		\item \label{enum:19:1}Für $^*$-Darstellungen $\pi_A \colon A \to \mathcal{B}(\mathcal{H}_1)$, $\pi_B \colon B \to \mathcal{B}(\mathcal{H}_2)$ definiert 
		\[
			(\pi_A \otimes \pi_B)(a \otimes b) := \iota\enbrace[\big]{\pi_A(a) \otimes \pi_B(b)}
		\]
		eine $^*$-Darstellung $\pi_A \otimes \pi_B \colon A \odot B \to \mathcal{B}(\mathcal{H}_1 \otimes \mathcal{H}_2)$.
		Falls $\pi_A$ und $\pi_B$ treu sind, so ist auch $\pi_A \otimes \pi_B$ treu.
		\item Die Abbildung $\norm*{\cdot }_{\min} \colon A \odot B \to \mathbb{R}^+$, gegeben durch
		\[
			\norm[\big]{\textstyle\sum\nolimits_{i=1}^{n} a_i \otimes b_i}_{\min} := \sup \set[\Big]{\norm[\big]{(\pi_A \otimes \pi_B)\enbrace[\big]{\textstyle\sum\nolimits_{i=1}^{n} a_i \otimes b_i}} \given \pi_A, \pi_B \text{ $^*$-Darstellungen} }
		\]
		ist eine $C^*$-Norm.
	\end{enumerate}
\end{proposition}
\begin{beweis}
	\begin{enumerate}[(i)]
		\item Nach \autoref{bem:14} und \autoref{bem:16} existiert genau eine $^*$-Darstellung 
		\[
			\pi_A \otimes \pi_B = \iota \circ (\pi_A \otimes \pi_B)
		\]
		Seien $\pi_A$ und $\pi_B$ treu und $0 \neq c \in \ker (\pi_A \otimes \pi_B)$.
		Es gilt $c= \sum_{i=1}^{n} a_i \otimes b_i$ für geeignete $a_i$, $b_i$ und $n$.
		Ohne Einschränkungen können wir annehmen, dass die $b_i$ linear unabhängig sind.
		Da $\pi_B$ treu ist, sind auch die $\pi_B(b_i)$ linear unabhängig.
		Für $\xi \in \mathcal{H}_1$ wählen eine Orthonormalbasis $\set*{\xi_1, \ldots ,\xi_m}$ für $\Span \set*{\pi_A(a_i) \xi \given i=1,\ldots ,n} \subset \mathcal{H}_1$.
		Es existieren also $\lambda_{ij} \in \mathbb{C}$ für $i=1,\ldots ,n$, $j=1,\ldots ,m$ mit 
		\(
			\pi_A(a_i) \xi = \sum_{j=1}^{m} \lambda_{ij} \cdot \xi_j
		\)
		für $i=1,\ldots ,n$.
		Für $\eta \in \mathcal{H}_2$ gilt
		\begin{align}
			0 = \enbrace*{\pi_A \otimes \pi_B}(c) \enbrace*{\xi \otimes \eta} = \sum_{i=1}^{n} \pi_A(a_i) \xi \otimes \pi_B(b_i) \eta &= \sum_{i=1}^{n} \enbrace*{\sum_{j=1}^{m} \lambda_{ij} \cdot \xi_j} \otimes \pi_B(b_i) \eta \\
			&= \sum_{j=1}^{m} \xi_j \otimes \enbrace*{\sum_{i=1}^{n} \lambda_{ij} \cdot \pi_B(b_i) \eta }
		\end{align}
		Es folgt, dass $\sum_{i=1}^{n} \lambda_{ij} \cdot \pi_B(b_i) \eta=0$ für $j=1,\ldots ,m$ und $\eta \in \mathcal{H}_2$.
		Damit sind die $\pi_B(b_i)$ linear abhängig oder $\lambda_{ij}=0$.
		Ersteres kann nicht gelten, also muss $\pi_A(a_i)\xi=0$ für $\xi \in \mathcal{H}_1$ gelten, also $\pi_A(a_i)=0$.
		Folglich ist $c=0$ im Widerspruch zur Annahme.
		\item $C^*$-Halbnorm: Übung!
		Wir zeigen, dass $\norm*{\cdot }_{\min}$ endlich ist: 
		\begin{align}
			\norm*{\enbrace*{\pi_A \otimes \pi_B} \enbrace*{\sum\nolimits_{i=1}^{n} a_i \otimes b_i}} &= \norm*{\sum_{i=1}^{n} \enbrace*{\pi_A(a_i) \otimes \ind_{\mathcal{H}_1} } \enbrace*{\ind_{\mathcal{H}_2} \otimes \pi_B(b_i)}}\\
			&\le \sum_{i=1}^{n} \norm*{\pi_A(a_i) \otimes \ind_{\mathcal{H}_1}} \norm*{\ind_{\mathcal{H}_2} \otimes \pi_B(b_i)} \\
			\intertext{$\pi_A(\cdot ) \otimes \ind_{\mathcal{H}_1}$ und $\ind_{\mathcal{H}_2} \otimes \pi_B(\cdot )$ sind als $^*$-Darstellungen kontraktiv, also}
			&\le \sum_{i=1}^{n} \norm*{a_i} \norm*{b_i} < \infty
		\end{align}
		$\norm*{\cdot }_{\min}$ ist eine Norm nach (i), denn $A$ und $B$ besitzen treue Darstellungen.\qedhere
	\end{enumerate}
\end{beweis}

\begin{definition}[{name=[{minimales Tensorprodukt}]}]
	\[
		\overline{A \odot B}^{\norm*{\cdot }_{\min}} =: A \otimes_{\min}B \enspace(= {A \otimes B})
	\]
	heißt \bet{minimales} (oder \bet{räumliches} oder \bet{spatiales}) \bet{Tensorprodukt}\index{minimales Tensorprodukt} von $A$ und $B$.
\end{definition}

\begin{satz}[{name=[{Zerlegung von Darstellungen des algebraischen Tensorproduktes}]},label=satz:111]
	Seien $A,B$ $C^*$-Algebren und $\pi \colon A \odot B \to \mathcal{B}(\mathcal{H})$ eine nichtdegenerierte Darstellung.
	Dann existieren eindeutig bestimmte nichtdegenerierte Darstellungen $\pi_A \colon A \to \mathcal{B}(\mathcal{H})$ und $\pi_B \colon B \to \mathcal{B}(\mathcal{H})$ mit 
	\[
		\pi(a \otimes b) = \pi_A(a) \cdot \pi_B(b)=\pi_B(b) \cdot \pi_A(a)
	\]
	Falls $\pi$ eine Faktordarstellung ist (das heißt $\overline{\pi(A \odot B)}^{\wo}$ ist ein Faktor), so auch $\pi_A$ und $\pi_B$.
\end{satz}
\begin{beweis}
	Falls $A$ und $B$ unital sind, so sind $\pi_A$, $\pi_B$ gegeben durch $\pi_A(a) = \pi(a \otimes \ind_B)$ bzw. $\pi_B(b) = \pi(\ind_A \otimes b)$.
	Im allgemeinen Fall seien $(h_\lambda), (k_\mu)$ approximative Einsen.
	Setze 
	\begin{align}
		\pi_A(a) &:= \solim \pi(a \otimes k_\mu) \\
		\pi_B(b) &:= \solim \pi(h_\lambda \otimes b)
	\end{align}
	Wir zeigen, dass diese Limiten existieren:
	Für $\xi \in \mathcal{H}$ und $a \in A_+$ definiert $b \mapsto \skal*{\xi}{\pi(a \otimes b) \xi}$ ein positives Funktional $f_a$ auf $B$ (beschränkt nach \cite[Satz 5.3]{OpAlg}).
	Damit existiert
	\[
		\lim_\mu f_a(k_\mu) = \lim_\mu \skal*{\xi}{\pi(a \otimes k_\mu)\xi}
	\]
	Für $a \in A_+$, $\mu \le \mu'$ gilt 
	\begin{align}
		\norm[\Big]{\enbrace[\big]{\pi(a \otimes k_{\mu'}) - \pi (a \otimes k_\mu)}\xi}^2 = \skal[\Big]{\xi}{\pi \enbrace[\big]{a^2 \otimes \enbrace*{k_{\mu'} - k_\mu}^2} \xi}
		&\le \skal[\Big]{\xi}{ \pi \enbrace[\big]{a \otimes \enbrace*{k_{\mu'}- k_\mu}} \xi} \\
		&= f_a(k_{\mu'}) - f_a(k_\mu) < \varepsilon
	\end{align}
	für $\mu, \mu'$ groß genug.
	Also existiert $\pi_A(a)$ für $a \in A_+$. 
	Mit Linearität folgt, dass $\pi_A$ auch insgesamt wohldefiniert ist.
	Dass $\pi_A, \pi_B$ $^*$-Darstellungen sind mit den gewünschten Eigenschaften, verbleibt als Übung!
	
	Es gilt nach dem Bikommutantensatz
	\[
		\begin{tikzcd}
			\pi_A(A)\dar[phantom,"\subset", sloped] \rar[phantom,"\subset", sloped] & \pi_B(B)' \dar[equals] \\
			\pi_A(A)'' \rar[phantom,"\subset", sloped] & \pi_B(B)''
		\end{tikzcd}
	\]
	Damit folgt $\pi_A(A)' \cap \pi_A(A)'' \subset \pi_A(A)' \cap \pi_B(B)' \subset \pi(A \odot B)'$.
	Weiter haben wir $\pi_A(A) \subset \overline{\pi(A \odot B)}^{\so} = \pi(A \odot B)''$, also gilt $\pi_A(A)'' \subset \pi(A \odot B)''$, womit weiter folgt
	\[
		\mathbb{C} \cdot \ind_{\mathcal{H}} \subset \pi_A(A)' \cap \pi_A(A)'' \subset \pi(A \odot B)' \cap \pi(A \odot B)'' = \mathbb{C} \cdot \ind_{\mathcal{H}}\qedhere
	\]
\end{beweis}

\begin{korollar}[{name=[{Submultiplikativität bezüglich des Tensorproduktes}]},label=kor:112]
	Seien $A,B$ $C^*$-Algebren, $\gamma$ eine $C^*$-Halbnorm auf $A \odot B$.
	Dann gilt $\gamma(a \otimes b) \le \norm*{a} \cdot \norm*{b}$ für $a \in A$, $b	\in B$.
	$\gamma$ besitzt eine eindeutige Fortsetzung zu einer $C^*$-Halbnorm auf $A^\simm \odot B^\simm$.
\end{korollar}
\begin{beweis}
	Setze $N := \set*{x \in A \odot B \given \gamma(x)=0} \lhd A \odot B$.
	$\gamma$ induziert eine $C^*$-Norm $\overline{\gamma}$ auf $\sfrac{A \odot B}{N}$.
	Damit ist $C := \overline{\sfrac{A \odot B}{N}}^{\overline{\gamma}}$ eine $C^*$-Algebra.
	$C$ besitzt eine nichtdegenerierte treue Darstellung.
	Wir erhalten $\pi \colon A \odot B \to \mathcal{B}(\mathcal{H})$ mit $\norm*{\pi(X)}= \overline{\gamma}(x +N) =\gamma(x)$ für $x \in A \odot B$.
	Weiter gilt
	\[
		\gamma(a \otimes b) = \norm*{\pi(a \otimes b)} \le \norm*{\pi_A(a)} \cdot \norm*{\pi_B(b)} \le \norm*{a} \cdot \norm*{b}
	\]
	Seien $\tilde{\pi}_A \colon A^\simm \to \mathcal{B}(\mathcal{H})$, $\tilde{\pi}_B \colon B^\simm \to \mathcal{B}(\mathcal{H})$ die Fortsetzungen auf den Unitalisierungen.
	Dann definiert $\tilde{\pi}(a \otimes b) := \tilde{\pi}_A(a) \cdot \tilde{\pi}_B(b)$ eine Fortsetzung $\tilde{\pi} \colon A^\simm \odot B^\simm \to \mathcal{B}(\mathcal{H})$ von $\pi$.
	($\tilde{\pi}$ existiert nach \autoref{bem:14}).
	$\tilde{\gamma}(\cdot ) := \norm*{\tilde{\pi}(\cdot)}$ ist eine Fortsetzung von $\gamma$.
	Wir müssen die Eindeutigkeit zeigen:
	\begin{align}
		\tilde{\gamma}(x) = \norm*{\tilde{\pi}(x)} = \norm*{\ind_{\mathcal{H}} \cdot \tilde{\pi}(x)} &= \norm*{\solim \pi\enbrace[\big]{h_\lambda^{\sfrac{1}{2}} \otimes k_\mu^{\sfrac{1}{2}}} \cdot \tilde{\pi}(x)} \\
		&= \lim_{\lambda,\mu} \norm[\big]{\tilde{\pi}(x^*) \pi \enbrace*{h_\lambda \otimes k_\mu} \tilde{\pi}(x)}^{\sfrac{1}{2}} \\
		&= \lim_{\lambda,\mu} \norm[\big]{\pi \enbrace*{x^* (h_\lambda \otimes k_\mu) x}}^{\sfrac{1}{2}} \\
		&= \lim_{\lambda,\mu} \gamma \enbrace[\big]{x^*(h_\lambda \otimes k_\mu)x}
	\end{align}
	Der letzte Ausdruck ist unabhängig von $\pi$.
\end{beweis}

\begin{definitionP}[{name=[{Max-Norm}]}]
	Die Abbildung $\norm*{\cdot }_{\max} \colon A \odot B \to \mathbb{R}^+$ gegeben durch
	\[
		\norm*{x}_{\max} := \sup \set[\big]{\norm*{\pi(x)} \given \pi \text{ ist $^*$-Darstellung von } A \odot B}
	\] 
	ist eine $C^*$-Norm auf $A \odot B$.
	Für jede weitere $C^*$-Norm $\gamma$ gilt $\gamma(x) \le \norm*{x}_{\max}$ für $x \in A \odot B$.
\end{definitionP}
\begin{beweis}
	$\norm*{\cdot}_{\max}$ ist submultiplikativ, also ist $\norm*{x}_{\max}$ endlich für alle $x \in A \odot B$.
	Damit ist dann auch klar, dass $\norm*{\cdot}_{\max}$ eine $C^*$-Halbnorm ist.
	Es seien nun $\pi_A, \pi_B$ treue Darstellungen von $A$ beziehungsweise $B$.
	Dann ist $\pi_A \otimes \pi_B$ eine treue Darstellung von $A \odot B$ nach \autoref{prop:19} \ref{enum:19:1}.
	Damit ist $\norm*{\cdot}_{\max}$ eine $C^*$-Norm und auch die Maximalität folgt.
\end{beweis}

\begin{definition}[{name=[{maximales Tensorprodukt}]}]
	$A \tensormax B := \overline{A \odot B}^{\norm*{\cdot }_{\max}}$ heißt \Index{maximales Tensorprodukt} von $A$ und $B$. 
\end{definition}

\begin{bemerkung}
	\begin{enumerate}[(i)]
		\item Sei $\gamma$ eine $C^*$-Norm auf $A \odot B$.
		Dann erhalten wir eine Abbildung
		\[
			\begin{tikzcd}
				A \tensormax B \rar[two heads] & A \otimes_\gamma B
			\end{tikzcd}
		\]
		denn $\id \colon A \odot B \to A \odot B$ ist normvermindernd, setzt sich also fort auf $A \tensormax B$.
		Das Bild ist eine dichte $C^*$-Unteralgebra, also alles.
		\item Sind $A,B$ unital, so ist $A \otimes_\gamma B$ unital und es gibt Inklusionen
		\begin{align}
			A &\xhookrightarrow{\hspace{1.4em}} A \mathop{\otimes_\gamma} B , \qquad a \longmapsto a \otimes \ind_B \\
			B &\xhookrightarrow{\hspace{1.4em}} A \mathop{\otimes_\gamma} B , \qquad b \longmapsto \ind_A \otimes b
		\end{align}
		\item Sind $A$, $B$ unital, dann ist $A \tensormax B \cong C^* \enbrace*{A,B \mid \benbrace*{A,B}=0}$ (universelle $C^*$-Algebra, Übung!)
		\item $\tensormax $ und $\tensormin$ sind assoziativ und kommutativ (warum?).
	\end{enumerate}
\end{bemerkung}

\begin{definition}[{name=[{Zustände auf dem algebraischen Tensorprodukt}]}]
	Für $C^*$-Algebren $A$, $B$ setzen wir 
	\[
		(A \odot B)^* := \set[\big]{f \colon A \odot B \to \mathbb{C} \text{ linear}}
	\]
	$f \in (A \odot B)^*$ heißt positiv, falls $f(x^* x)\ge 0$ für $x \in A \odot B$.
	Wir setzen
	\[
		\norm*{f} := \sup \set*{f(a \otimes b) \given a \in A^1_+, b \in B_+^1}
	\]
	Wir definieren die \bet{Zustände}\index{Zustände auf dem algebraischen Tensorprodukt} auf $A \odot B$ durch $S(A \odot B) := \set[\big]{f \in (A \odot B)^* \given \norm*{f}=1}$.
\end{definition}

\begin{proposition}
	Jedes positive $f \in (A \odot B)^*$ besitzt eine positive Fortsetzung $f^\sim \in \enbrace*{A^\simm \odot B^\simm}^*$ mit 
	\[
		\norm*{f}= \norm*{f^\sim} = \tilde{f}(\ind_{A^\simm} \otimes \ind_{B^\simm}) < \infty
	\]
\end{proposition}
\begin{beweis}
	Sei $(h_\lambda)_{\Lambda} \subset A$ eine approximative Eins. 
	Für $b \in B_+$ setze $f_b(a) := f (a \otimes b)$.
	Dann ist $f_b \in A^*$ positiv und insbesondere beschränkt.\todo{Ref}
	Wir setzen
	\[
		\hat{f} \enbrace*{\ind_{A^\simm} \otimes b} := \Underbracket{\lim_\lambda f(h_\lambda \otimes b)}{f^\sim_b (\ind_{A^\simm})} = \sup_\lambda f(h_\lambda \otimes b) < \infty
	\]
	Wir setzen $\hat{f}$ wie folgt fort, sodass $\hat{f} \in (A^\simm \odot B)^*$ gilt:
	\[
		\hat{f}(x) = \lim_\lambda f \enbrace[\big]{\enbrace*{h_\lambda \otimes  \ind_{B^\simm}} x} = \lim_\lambda f\enbrace*{\enbrace*{h_\lambda \otimes  \ind_{B^\simm}}^{\sfrac{1}{2}} x \enbrace*{h_\lambda \otimes  \ind_{B^\simm}}^{\sfrac{1}{2}}}
	\]
	Man sieht damit, dass $\hat{f}$ linear, wohldefiniert und positiv ist. 
	$\hat{f}$ lässt sich fortsetzen zu $f^\sim \in (A^\simm \odot B^\simm)^*$ positiv.
	Es gilt trivialerweise $\norm*{f} \le \norm*{f^\sim}$ und $f^\sim (a \otimes b) \le f^\sim (\ind_{A^\simm} \otimes \ind_{B^\simm})$ für $a \in (A_+^\simm)^1$ und $b \in (B_+^\simm)^1$.
	Damit folgt $\norm*{f^\sim} \le f^\sim (\ind_{A^\simm} \otimes \ind_{B^\simm}) \le \norm*{f^\sim}$.
	Weiter gilt
	\[
		f^\sim \enbrace*{\ind_{A^\simm} \otimes \ind_{B^\simm}} = \lim_{\lambda,\mu} f \enbrace*{h_\lambda \otimes k_\mu} \le \norm*{f}
	\]
	Damit ist $\norm*{f} = \norm*{f^\sim} = f^\sim \enbrace*{\ind_{A^\simm} \otimes \ind_{B^\simm}}$ wie behauptet.
\end{beweis}

\begin{lemma}[label=lem:118]
	Seien $A$, $B$ unitale $C^*$-Algebren und $x=x^* \in A \odot B$.
	Dann existieren $\lambda > 0$, $n \in \mathbb{N}$, $z_1, \ldots ,z_n \in A \odot B$ mit 
	\[
		\lambda \cdot \enbrace*{\ind_{A^\simm} \otimes \ind_{B^\simm}} -x = \sum_{i=1}^{n} z_i^* \, z_i
	\]
	Wir schreiben auch $\lambda \cdot \enbrace*{\ind_{A^\simm} \otimes \ind_{B^\simm}} \ge_{\alg} x$.
\end{lemma}
\begin{beweis}
	Für $a \in A_+$, $b \in B_+$ gilt
	\[
		0 \le_\alg a \otimes b \le_\alg \norm*{a} \cdot \norm*{b} \cdot \enbrace*{\ind_A \otimes \ind_B}
	\]
	Weiter gilt $(A \odot B)_\so = A_\so \odot B_\so$ \emph{nicht}.
	Falls $y=y* = \sum_{k=1}^{m} a_k \otimes b_k$, so gilt
	\[
		y= \frac{1}{4} \cdot \sum_{k=1}^{m} (a_k +a_k^*) \otimes (b_k + b_k^*) - i (a_k - a_k^*) \otimes i \cdot (b_k -b_k^*)
	\]
	Schreibe $x= \sum_{k=1}^{n} c_k \otimes d_k$ mit $c_k \in A_\sa$, $d_k \in B_\sa$.
	Dann ist 
	\[
		x \le_\alg \sum_{k=1}^{n} \enbrace[\big]{(c_k)_+ + (c_k)_-} \otimes  \enbrace[\big]{(d_k)_+ + (d_k)_-} \le_\alg \enbrace[\bigg]{\Underbracket{\sum \norm*{c_k} \cdot \norm*{d_k}}{=: \lambda} } \cdot \enbrace*{\ind_A \otimes \ind_B} \qedhere
	\]
\end{beweis}

Sei $f \in (A \odot B)^*$ positiv. 
Die GNS-Konstruktion liefert einen Hilbertraum 
\[
	\mathcal{H}_f := \overline{\sfrac{A \odot B}{N_f}}^{\skal*{\cdot}{\cdot }_f}
\]
Mit Vektoren $\xi_f = \benbrace*{\ind_A \otimes \ind_B} \in \mathcal{H}_f$.
Falls $A$ oder $B$ nicht unital sind, benutze $\tilde{f} \in (A^\simm \odot B^\simm)^*$.
Wir definieren einen $^*$-Homomorphismus $\pi_f \colon A \odot B \to \mathcal{B}(\mathcal{H}_f)$ wie folgt:
Für $x,y \in A \odot B$ setze $\pi_f(x)[y] := [xy] \in \sfrac{A \odot B}{N_f} \subset \mathcal{H}_f$.
$\pi_f(x)$ setzt sich auf $\mathcal{H}_f$ fort: 
Mit \autoref{lem:118} folgt die Existenz von $\lambda>0$ mit $x^*x \le_\alg \lambda \cdot (\ind_{A} \otimes \ind_B)$.
Für $y \in A \odot B$ gilt 
\begin{align}
	\norm*{\pi_f(x)[y]}^2 = \norm*{[xy]}^2 = f \enbrace*{y^* x^*x y} &= f \enbrace*{y^* \lambda  \cdot (\ind_{A} \otimes \ind_B) y} - \sum_{k=1}^{n} f \enbrace*{y^* z_k^* z_k y}\\
	&\le \lambda \cdot 	f(y^* y) = \lambda \cdot \norm*{[y]}^2
\end{align}
Also ist $\norm*{\pi_f(x)} \le \lambda^{\sfrac{1}{2}}$ und damit ist $\pi_f(x)$ beschränkt und setzt sich fort auf $\mathcal{H}_f$.
$\pi_f$ induziert einen $^*$-Homomorphismus $\overline{\pi}_f \colon A \tensormax B \to \mathcal{B}(\mathcal{H}_f)$.
Dieser wiederum induziert eine eindeutige stetige Fortsetzung 
\[
	\overline{f}( .) := \skal*{\xi_f}{\overline{\pi}_f(.) \xi_f}
\]
von $f$ auf $A \tensormax B$.
Wir haben nun einen Homöomorphismus $S(A \odot B) \xrightarrow{\approx} S(A \tensormax B)$ konstruiert.
Wir erhalten:

\begin{satz}[label=satz:119]
	Für $x \in A \odot B$ gilt
	\[
		\norm*{x}_{\max}^2 = \sup \set*{\frac{f(y^* x^* xy)}{f(y^* y)} \given f \in S(A \odot B) , y \in A \odot B, f(y^* y)\neq 0 }
	\]
\end{satz}
\begin{beweis}
	Es ist
	\begin{align}
		&\hspace{1.14em}\sup \set*{\frac{f(y^* x^* xy)}{f(y^* y)} \given f \in S(A \odot B) , y \in A \odot B, f(y^* y)\neq 0 }\\ 
		&=\sup \set*{\frac{f(y^* x^* xy)}{f(y^* y)} \given f \in S(A \tensormax B) , y \in A \odot B, f(y^* y)\neq 0 } \\
		&= \sup \set*{\norm*{\pi_f(x)}^2 \given f \in S(A \tensormax B)} \\
		&= \norm*{x}_{\max}^2 \qedhere
	\end{align}
\end{beweis}

Wir wollen nun ein ähnliches Ergebnis für die Minimumsnorm beweisen:
Seien $f \in A^*$, $g \in B^*$ positiv, dann ist $f \otimes g \in (A \odot B)^*$ positiv, wo $(f \otimes g) \enbrace*{\sum_{i=1}^{n} a_i \otimes b_i} = \sum_{i=1}^n f(a_i) g(b_i)$.
Seien $(\pi_f,\xi_f)$, $(\pi_g,\xi_g)$ die GNS-Darstellungen von $A$ beziehungsweise $B$.
Dann ist
\[
	(f \otimes g)(x) = \skal[\big]{ \enbrace*{\xi_f \otimes \xi_g} }{(\pi_f \otimes \pi_g)(x) \enbrace*{\xi_f \otimes \xi_g}}
\]
für $x \in A \odot B$.
$f \otimes g$ ist also ein Vektorzustand und damit positiv.
Wir erhalten:

\begin{satz}[label=satz:120]
	Es gilt für $x \in A \odot B$
	\[
		\norm*{x}^2_{\min} = \sup \set*{\frac{(f \otimes g) (y^*x^*xy)}{(f \otimes g)(y^*y)} \given f \in S(A), g \in S(B), y \in A \odot B, (f \otimes g)(y^* y)\neq 0 }
	\]
	(Tatsächlich genügt es reine Zustände zu betrachten.)
\end{satz}
\begin{beweis}
	Es gilt für $x \in A \odot B$\marginnote{nach Zorn ist jede Darstellung Summe zyklischer Darstellungen}
	\begin{align}
		\norm*{x}_{\min}^2 &= \sup \set*{\norm*{(\pi_A \otimes \pi_B)(x)}^2 \given \pi_A, \pi_B \text{ Darstell.} } \\
		&= \sup \set*{ \norm*{(\pi_A \otimes \pi_B)(x)}^2 \given (\pi_A,\xi_A), (\pi_B,\xi_B) \text{ zykl. Darstell.}} \\
		&= \sup \set*{ \frac{\skal[\Big]{(\pi_A \otimes \pi_B)(y) \enbrace*{\xi_A \otimes \xi_B}}{(\pi_A \otimes \pi_B)(x^*x) \enbrace[\big]{ (\pi_A \otimes \pi_B)(y) \enbrace*{\xi_A \otimes \xi_B}}}}{\norm[\big]{(\pi_A \otimes \pi_B)(y) \enbrace*{\xi_A \otimes \xi_B}}^2}  \given\hspace{-3pt} \begin{array}{l}
	(\pi_A,\xi_A), \\
	(\pi_B, \xi_B)
\end{array} \hspace{-3pt}} \\
		&= \sup \set*{ \frac{\enbrace*{f_{\xi_A} \otimes f_{\xi_B}} \enbrace*{y^*x^*xy}}{\enbrace*{f_{\xi_A} \otimes f_{\xi_B} }(y^*y)} \given (\pi_A,\xi_A), (\pi_B, \xi_B)} \\
		&\le \sup \set*{ \frac{(f \otimes g)(y^*x^* xy)}{(f \otimes g)(y^*y)} \given f \in S(A) , g \in S(B), y \in A \odot B, (f \otimes g)(y^* y)\neq 0 } \\
		&\StackTextClap{GNS}{\le} \norm*{x}_{\min}^2 \qedhere
	\end{align}
\end{beweis}

\begin{satz}[{name={Takesaki}},label=satz:121]
	Seien $A =C_0(X)$ und $B$ $C^*$-Algebren.
	Dann existiert nur eine $C^*$-Norm auf $A \odot B$ und $A \otimes B \cong C_0(X,B)$ (mit $a \otimes b(x) = a(x) \cdot b$).
	Falls $B=C_0(Y)$, so haben wir $A \otimes B \cong C_0(X \times Y)$ mit $(a \otimes b)(x,y)=a(x) \cdot b(y)$.
\end{satz}
\begin{beweis}
	In mehreren Schritten
	\begin{enumerate}[(i)]
		\item Wir dürfen annehmen, dass $A$ und $B$ unital sind: 
		\[
			\norm*{\cdot }_{\min, A^\simm \odot B^\simm}\big|_{A \odot B} = \norm*{\cdot }_{\min, A \odot B}
		\]
		denn $\pi_A^\simm \otimes \pi_B^\simm\big|_{A \odot B} = \pi_A \otimes \pi_B$.
		Ebenso für $\norm*{\cdot }_\gamma$, denn $\gamma$ setzt sich fort zu $\gamma^\simm$ auf $A^\simm \otimes B^\simm$ (nach \autoref{satz:111} bzw. \autoref{kor:112}).
		\item Es gilt $P(A \tensormax B) \approx P(A) \times P(B)$:
		
		Sei $f \in P(A \tensormax B)$.
		Die GNS-Darstellung $\pi_f$ ist dann irreduzibel.\todo{Ref letztes Semester}
		Nach \autoref{satz:111} erhalten wir Darstellungen $\pi_A$, $\pi_B$ von $A$ bzw. $B$.
		$\pi_f$ ist eine Faktordarstellung, also ist $\pi_A$ nach \autoref{satz:111} auch eine Faktordarstellung.
		Nun gilt, da $A$ kommutativ ist
		\[
			\pi_A(A) \subset \pi_A(A)'' \cap \pi_A(A)'= C \cdot \ind_{\mathcal{H}_f} \implies \pi_A \in \hat{A}
 		\]
		also ist $\pi_A(\cdot ) = \ev_x(\cdot) \cdot \ind_{\mathcal{H}_f}$ für ein $x \in X$.
		Damit ist 
		\[
			f(a \otimes b) = \skal[\big]{\xi_f}{\pi_f(a \otimes b) \xi_f} = \ev_x(a) \cdot \Underbracket{\skal[\big]{\xi_f}{\pi_f(\ind_A \otimes b) \xi_f}}{=: g(b)}
		\]
		$\pi_B$ ist irreduzibel und $g$ ist Vektorstand und damit rein, also $g \in P(B)$ und damit $f= \ev_x \otimes g$.
		Für die Umkehrung: $\ev_x \otimes h \in P(A \tensormax B)$ (Warum?)
		\item Wir zeigen $\norm*{\cdot }_{\max} = \norm*{\cdot }_{\min}$:
		\[
			\hspace{-1em}\norm*{z}^2_{\max} = \sup \set[\big]{f(z^*z) \given f \in P(A \tensormax B)} \le \sup \set[\big]{(\ev_x \otimes g)(z^*z) \given x \in X, g \in P(B)} \StackTextClap{\ref{satz:120}}{\le} \norm*{z}_{\min}^2 \le \norm*{z}_{\max}^2
		\]
		für $z \in A \odot B$.
		\item Es gilt $\Prim(A \tensormax B) \approx X \times \Prim B$:\todo{hübsch machen}
		\[
			\begin{tikzcd}
				& {\ev_x} \otimes g & (x,g) \lar[mapsto] \\
				{\ev_x} \otimes g \dar[mapsto] & P(A \tensormax B) \dar[two heads]& X\times P(B) \lar \dar[two heads] & (x,g) \dar[mapsto] \\
				\ker(\pi_{\ev_x \otimes g}) = \ker(\pi_{\ev_x} \otimes \pi_g) & \Prim(A \tensormax B) & X \times \Prim B \lar[dotted,"\Theta"] & (x, \ker \pi_g) \\
				& \ker (\pi_{\ev_x} \otimes \pi_g) & (x, \ker \pi_g) \lar[mapsto]
			\end{tikzcd}
		\]
		$\Theta$ ist wohldefiniert (dann ist $\Theta$ automatisch ein Homöomorphismus):
		\[
			\begin{tikzcd}[column sep=0pt,row sep=2.5em]
				& & A \tensormax B \dlar["\delta_x"] \ar[ddrr,"\ev_x \otimes \pi_g = \pi_{\ev_x \otimes g}"]\\
				& B \dlar["\pi_g"] \\
				\mathcal{B}(\mathcal{H}_g) \ar[rrrr,"\cong"] & & & & \mathcal{B} \enbrace*{C \otimes \mathcal{H}_g}
			\end{tikzcd}
		\]
		Damit ist $\ker \pi_{\ev_x \otimes g} = \delta^{-1}_x (\ker \pi_g)$.
		Dies hängt nur von $\ker \pi_g$ ab und damit ist $\Theta$ wohldefiniert.
		\item $\norm*{\cdot}_{\max}=\norm*{\cdot}_{\min}$ ist die einzige $C^*$-Norm auf $A \odot B$:
        
        Sei $\gamma$ eine weitere $C^*$-Norm auf $A \odot B$ und $q \colon A \tensormax B \twoheadrightarrow  A \otimes_\gamma B$ die Quotientenabbildung.
        Dann ist 
        \[
            Y \coloneqq \set*{y \in \Prim(A \tensormax B) \given \ker q \subset y} \subset_{\mathrm{abg.}} \Prim \enbrace*{A \tensormax B}
        \]
        Erinnerung an die Hülle-Kern-Topologie:\todo{hübsch machen} 
        \[
            \begin{tikzcd}
                \set*{\text{abg. Ideale von }C} \rar["\approx"] & \set*{\text{abg. Teilmengen von }\Prim C } \\
                J \rar[mapsto] & \set*{y \in \Prim C \given J \subset y} \\
                \bigcap_{y \in W} y & W \lar[mapsto]
            \end{tikzcd}
        \]
        Angenommen $\set*{0} \neq \ker q$.
        Dann ist $Y \neq \Prim (A \tensormax B) \approx X \times \Prim B$.
        Folglich existiert $\emptyset \neq U \subset \Prim A =X$ offen und $\emptyset \neq V \subset \Prim B$ offen, sodass $U \times V \cap Y = \emptyset$.
        Dann sind $X \setminus U \subsetneq X$ und $\Prim B \setminus V \subsetneq \Prim B$ abgeschlossen, sowie $I \coloneqq \bigcap_{x \in X \setminus U} \ker \ev_x \neq \set*{0}$.
        Weiter ist $I \lhd A$ ein Ideal.
        Auch
        \(
            J \coloneqq \bigcap_{z \in \Prim B \setminus V} z \neq \set*{0}
        \)
        ist ein Ideal in $B$.
        Wir haben nun\marginnote{nach (iv) kann man $y=(y_A,y_B)$ schreiben}
        \[
            \ker q = \bigcap_{y \in Y} y \supset \bigcap_{y \notin U \times V} y = \enbrace*{\bigcap_{y_A \notin U} y} \cap \enbrace*{\bigcap_{y_B \notin V} y} \supset I \odot B \cap A \odot J \supset I \odot J \neq \set*{0}
        \]
        Damit folgt $\gamma(I \odot J)=0$.
        Aber $\gamma|_{I \odot J}$ ist eine $C^*$-Norm, darf also nicht verschwinden. Widerspruch!\qedhere
	\end{enumerate}
\end{beweis}

\begin{satz}[{name=[Die minimale Norm ist minimal]}]
    Seien $A,B$ $C^*$-Algebren und $\gamma$ eine $C^*$-Norm auf $A \odot B$.
    Dann gilt $\norm*{x}_{\min} \le \gamma(x)$ für alle $x \in A \odot B$.
\end{satz}
\begin{beweis}
    Wir beweisen den Satz in mehreren Schritten:
    \begin{itemize}
        \item Wir dürfen $A$ und $B$ als unital annehmen:
        
        Nach \autoref{kor:112} setzt sich $\gamma$ auf $A^\simm \odot B^\simm$ fort.
        Weiter gilt $\norm*{\cdot}_{\min,A^\simm \odot B^\simm}\big|_{A \odot B} = \norm*{\cdot}_{\min, A \odot B}$. (warum?)\marginnote{Darstellungen fortsetzen}
        \item Für reine Zustände $f \in P(A)$ und $g \in P(B)$ ist $f \otimes g$ stetig bezüglich $\gamma$ und setzt zu einem Zustand auf $A \otimes_\gamma B$ fort:
        
        Setze $S_\gamma \coloneqq \set[\big]{(f,g) \in P(A) \times P(B) \given f \otimes g \text{ hat Fortsetzung zu Zustand auf } A \otimes_\gamma B}$.
        Dann ist $S_\gamma \subset P(A) \times P(B)$ abgeschlossen in der Produkttopologie der $\w^*$-Topologien (Übung).
        Für $b \in B_+$ setze $C \coloneqq C^*(b) \cong C_0 \enbrace*{\sigma(b) \setminus \set*{0}}$.
        Wähle $\omega \in P(C)$ mit $\omega(b)=\norm*{b}$.\todo{Ref letztes Semester}
        Falls $f' \in P(A)$, dann ist $f' \otimes \omega \in P(A \tensormax C)= P (A \otimes_\gamma C)$ wie im Beweis von \autoref{satz:121} gesehen.
        Nach Krein-Milmann existiert eine Fortsetzung $\Theta \in P(A \otimes_\gamma B)$.\todo{ref  7.14 letztes Semester}
        
        Es gilt $\Theta = f' \otimes g'$, wobei $g' \coloneqq \Theta|_{\ind_A \otimes_\gamma B} \in P(B)$:
        Sei $y \in B_+^1$.
        Dann definieren wir $f_1 \in A^*$ durch $f_1(x) := \Theta(x \otimes y)$ und $f_2 \in  A^*$ durch $f_2(x) \coloneqq \Theta(x \otimes \ind_B -y)$.
        Wir haben dann $f_1 + f_2 =f'$ und da $f'$ rein ist, folgt $f_1 = \lambda \cdot f'$ mit $\lambda = f_1(\ind_A) =g'(y)$.
        Also ist
        \[
            \Theta(x \otimes y) =f_1(x) = f'(x) g'(y) = \enbrace*{f' \otimes g'}(x \otimes y) \implies \Theta = f' \otimes g'
        \]
        $g'$ ist ein reiner Zustand (warum?).
        Weiter gilt 
        \[
            \norm*{b} =\omega(b) =f' \otimes \omega(\ind_A \otimes b) = \Theta(\ind _A \otimes b) = (f' \otimes g')(\ind_A \otimes b) = g'(b)
        \]
        Insgesamt: $Y_{f'} \coloneqq \set*{g' \in P(B) \given (f',g') \in S_\gamma}$ normiert $B$ für jedes $f' \in P(A)$.
        Damit ist $Y_{f'} \subset P(B)$ $\w^*$-dicht (Übung).
        Aber $\set*{f'} \times Y_{f'} \subset S_\gamma \subset_{\mathrm{abg.}} P(A) \times P(B)$, also 
        \[
            \set*{f'} \times P(B) = \set*{f'} \times \overline{Y_{f'}}^{P(B)} = \overline{\set*{f'} \times Y_{f'}}^{P(A) \times P(B)} \subset S_\gamma
        \]
        Es folgt also $P(A) \times P(B) \subset S_\gamma \subset P(A) \times P(B)$.
        \item Jedes $f \otimes g$ mit $f \in S(A)$, $g \in S(B)$ setzt sich fort zu einem Zustand auf $A \otimes_\gamma B$ (Übung).
        \item Für $x \in A \odot B$ gilt
        \begin{align}
            \norm*{x}_{\min}^2 &\StackTextClap{\ref{satz:120}}{=} \sup \set*{ \frac{f \otimes g(y^* x^* xy)}{f \otimes g(y* y)} \given f \in S(A), g \in S(B), y \in A \odot B, \text{ Nenner } \neq 0} \\
            &\StackTextClap{(iii)}{\le} \sup \set*{ \frac{\gamma(x^* x) f \otimes g(y^*y)}{f \otimes g(y* y)} \given f \in S(A), g \in S(B), y \in A \odot B, \text{ Nenner } \neq 0}\\
            &= \gamma(x^*x) = \gamma(x)^2 \qedhere
        \end{align}
    \end{itemize}
\end{beweis}

\begin{bemerkung}
    \begin{enumerate}[(i)]
        \item Für jede $C^*$-Norm $\gamma$ auf $A \odot B$ gilt 
        \[
            \norm*{\cdot}_{\min} \le \gamma(\cdot ) \le \norm*{\cdot}_{\max}
        \]
        und $\gamma(a \otimes b) = \norm*{a}_A \cdot \norm*{b}_B$ für $a \in A$ und $b \in B$, denn
        \begin{align}
            \norm*{a}^2 \cdot \norm*{b}^2 \ge \gamma(a \otimes b)^2 \ge \norm*{a \otimes b}_{\min}^2 &\StackTextClap{\ref{satz:120}}{\ge} \sup \set[\big]{(f \otimes g) \enbrace*{a^* a \otimes b^*b} \given f \in P(A), g \in P(B)} \\
            &= \sup \set[\big]{f (a^*a) g(b^*b) \given f \in P(A), g \in P(B)} \\
            &= \norm*{a}^2 \cdot \norm*{b}^2
        \end{align}
        \item Wir haben kanonische $^*$-Homomorphismen 
        $A \tensormax B \twoheadrightarrow A \otimes_\gamma B \twoheadrightarrow A \tensormin B$
    \end{enumerate}
\end{bemerkung}

\begin{korollar}[{name=[minimales Tensorprodukt erhält Einfachheit]}]
    Sind $A$ und $B$ einfach, so ist $A \tensormin B$ auch einfach.
\end{korollar}
\begin{beweis}
	Es genügt zu zeigen, dass jede irreduzible Darstellung $\pi \colon A \tensormin B \to \mathcal{B}(\mathcal{H})$ isometrisch ist, denn wenn man für $J \lhd A \tensormin B$ die Komposition
	\[
		A \tensormin B \longrightarrow \sfrac{A \tensormin B}{J} \xrightarrow{\text{irred. Darst.}} \mathcal{B}(\mathcal{H})
	\]
	betrachtet, folgt dann $J=0$.
	Es genügt weiter zu zeigen, dass für $\pi \colon A \tensormin B \to \mathcal{B}(\mathcal{H})$ irreduzibel $\pi|_{A \odot B}$ treu ist, denn dann ist $\gamma(\cdot ) = \norm*{\pi(\cdot )}$ eine $C^*$-Norm auf $A \odot B$ und es gilt $\norm*{\cdot}_{\min} \le \gamma(\cdot) \le \norm*{\cdot }_{\min} \implies \pi$ isometrisch.
	
	Sei also $\pi$ irreduzibel gegeben und seien $\pi_A$, $\pi_B$ wie in \autoref{satz:111}.
	Sei $c = \sum_{i=1}^{n} a_i \otimes b_i \in A \odot B$ mit $\pi(c) = \sum_{i=1}^{n} \pi_A(a_i) \pi_B(b_i)=0$.
	Wir zeigen, dass $c=0$ ist.
	Setze $\tilde{\mathcal{H}} \coloneqq \mathcal{H} \otimes \mathbb{C}^n$ und
	\begin{align}
		X &\coloneqq \set*{\sum\nolimits_{i=1}^{n} \pi_B(bb_i) \xi \otimes e_i \given b \in B, \xi \in \mathcal{H}} \stackrel[\text{abg.}]{}{\subset} \tilde{\mathcal{H}} \\
		Y &\coloneqq \set*{\sum\nolimits_{i=1}^{n} \pi_A(a_i^*) \eta \otimes e_i \given \eta \in \mathcal{H} }\stackrel[\text{abg.}]{}{\subset} \tilde{\mathcal{H}}
	\end{align}
	Da $\pi(c) =0$ und $\benbrace*{\pi_A(A), \pi_B(B)}=0$ ist, gilt $X \perp Y$.
	Sei $\mathcal{B}(\tilde{\mathcal{H}}) \ni p \colon \tilde{\mathcal{H}} \to X \subset \tilde{\mathcal{H}}$ die orthogonale Projektion.
	Mit $\mathcal{B} (\tilde{\mathcal{H}}) \cong M_n \enbrace*{\mathcal{B}(\mathcal{H})}$ können wir $p=(p_{ij})$ mit $p_{ij} \in \mathcal{B}(\mathcal{H})$ schreiben.
	Es gilt $\enbrace*{\pi_B(B) \otimes \ind_{\mathbb{C}^n}} X \subset X$ und $\enbrace*{\pi_A(A) \otimes \ind_{\mathbb{C}^n}}X \subset X$, also ist
	\[
		p \in \enbrace[\big]{\pi_B(B) \otimes \ind_{\mathbb{C}^n}}' \cap \enbrace[\big]{\pi_A(A) \otimes \ind_{\mathbb{C}^n}}' \subset \mathcal{B}(\tilde{\mathcal{H}})
	\]
	Damit gilt dann auch $p_{ij} \in \pi_B(B)' \cap \pi_A(A)' \subset \mathcal{B}(\mathcal{H})$.
	Es folgt nun, dass $p_{ij} \in \pi \enbrace*{A \tensormin B}'$ gilt und weiter mit der Irreduzibilität $p_{ij} \in \mathbb{C}\cdot \ind_{\mathcal{H}}$, also $p \in \ind_{\mathcal{H}} \otimes M_n$.
	
	Für $b \in B$ und $\xi \in \mathcal{H}$ haben wir nun
	\begin{align}
		X \ni \sum_{i=1}^{n} \enbrace[\big]{\pi_B(bb_i) \xi} \otimes e_i = p \enbrace*{\sum_{i=1}^{n} \pi_B(bb_i) \xi \otimes e_i} &= \sum_{i,j=1}^n p_{ij} \pi_B(bb_j) \otimes e_i \\
		&= \sum_{i=1}^{n} \sum_{j=1}^{n} p_{ij} \pi_B(bb_j) \xi \otimes e_i
	\end{align}
	Also ist $\pi_B(bb_i)\xi = \sum_{i=1}^{n} p_{ij} \pi_B(bb_j) \xi$ für alle $b \in B$ und $\xi \in \mathcal{H}$.
	Da $b \in B$ beliebig war, folgt $\pi_B(b_i) = \sum_{j=1}^{n}  p_{ij} \pi_B(b_j)$ für $i=1,\ldots ,n$.
	Weiter gilt für $\eta \in \mathcal{H}$
	\[
		0 = p \Underbracket{\sum_{j=1}^{n} \pi_A(a_i^*) \eta \otimes e_i}{\in Y} = \sum_{i,j=1}^{n} p_{ij} \pi_B(a_j^*) \eta \otimes e_i
	\]
	Damit ist $0 = \sum_{j=1}^{n} p_{ij} \pi_A(a_j^*)$ für $i=1,\ldots ,n$, also auch $0 = \sum_{j=1}^{n} \overline{p_{ij}} \pi_A(a_j)$ für $i=1,\ldots ,n$.
	Wir erhalten also insgesamt\todo{Indizes noch falsch anscheinend}
	\[
		\pi_A \otimes \pi_B(c) = \sum_{i=1}^{n} \pi_A(a_i) \otimes \pi_B(b_i) = \sum_{i=1}^{n} \pi_A(a_i) \otimes  \sum_{j=1}^{n} p_{ij} \pi_B(b_j) = \sum_{i,j=1}^{n} p_{ij} \pi_A(a_i) \otimes \pi_B(b_j) =0
	\]
	$\pi_A$ und $\pi_B$ sind treu, also ist auch $\pi_A \otimes \pi_B$ treu auf $A \odot B$ und es folgt $c=0$.
\end{beweis}

\begin{proposition}[label=prop:125]
	Es seien $A_1$, $A_2$, $B_1$ und $B_2$ $C^*$-Algebren, $\varphi \colon A_1 \to B_1$ und $\psi \colon A_2 \to B_2$ $^*$-Homomorphismen.
	Dann existieren kanonische $^*$-Homomorphismen $\varphi \otimes \psi$
	\[
		\begin{tikzcd}[row sep=0pt]
			\varphi \otimes \psi \colon & A_1 \odot A_2 \rar & B_1 \odot B_2 \\
			\varphi \tensormax  \psi \colon & A_1 \tensormax  A_2 \rar & B_1 \tensormax  B_2 \\
			\varphi \tensormin  \psi \colon & A_1 \tensormin  A_2 \rar & B_1 \tensormin  B_2
		\end{tikzcd}
	\]
	Sind $\varphi$ und $\psi$ injektiv, so ist auch $\varphi \tensormin \psi$ injektiv -- im Allgemeinen ist dies nicht richtig für $\tensormax$.
\end{proposition}
\begin{beweis}
	\emph{Übung!}
\end{beweis}
% section 1 (end)
\newpage

\section{Nuklearität und Exaktheit} % (fold)
\label{sec:2}

\emph{\textenglish{Nucular. It's pronounced, \enquote{nucular}.}} -- Homer \textsc{Simpson}

\begin{definition}[{name=[{nukleare C*-Algebra}]}]
	Eine $C^*$-Algebra $A$ heißt \Index{nuklear}, falls für jede $C^*$-Algebra $B$ nur eine $C^*$-Norm auf $A \odot B$ existiert, das heißt $A \tensormax B = A \tensormin B$. 
\end{definition}

\begin{proposition}[{name=[{Nuklearität und Summen, Tensorprodukte, direkte Limiten und Quotienten}]},label=prop:22]
	\leavevmode
	\begin{enumerate}[(i)]
		\item Sind $A,B$ nuklear, so sind auch $A \oplus B$ und $A \otimes B$ nuklear.
		\item Ist $A = \Dlim\, A_i$ mit $A_i$ nuklear, so ist auch $A$ nuklear.
		\item Für \(
			\begin{tikzcd}[cramped,sep=small]
				0 \rar & J \rar & A \rar & \sfrac{A}{J} \rar & 0 
			\end{tikzcd}
		\) ist $A$ nuklear genau dann, wenn $J$ und $\sfrac{A}{J}$ nuklear sind.
	\end{enumerate}
\end{proposition}
\begin{beweis}[{name={mit \autoref{prop:24}}}]
	\leavevmode
	\begin{enumerate}[(i)]
		\item Für direkte Summen, ist die Aussage klar.
		Für Tensorprodukte auch wegen der Assoziativität von $\tensormax$ und $\tensormin$.\todo{RevChap 2}
		\item \enquote{klar} für injektive Verbindungsabbildungen. 
		Für den allgemeinen Fall benutzte (iii).
		\item Für die erste Implikation zeigen wir, dass $J$ nuklear ist.
		Folgendes Diagramm kommutiert
		\[
			\begin{tikzcd}
				J \tensormax  B \rar["\rho",two heads] & \overline{J \odot B}^{\norm*{\cdot}_{\max, A \odot B}} \dar[hook,two heads] \rar[phantom,sloped,"\lhd"] & A \tensormax \dar[equals] \\
				& J \tensormin B \rar[hook,"\ref{prop:125}"] & A \tensormin B
			\end{tikzcd}
		\]
		Damit ist $\overline{J \odot B}^{\norm*{\cdot}_{\max,A \odot B}} = J \tensormin B$.  
		Falls $\pi \colon J \odot B \to \mathcal{B}(\mathcal{H})$ eine nichtdegenerierte Darstellung ist und $(h_\lambda)$, $(g_\nu)$ approximative Einsen von $J$ bzw. $B$, so definieren wir $\sigma \colon A \odot B \to \mathcal{B}(\mathcal{H})$ durch
		\[
			\sigma(a \otimes b) \coloneqq \lim_{\lambda,\nu} \pi \enbrace[\big]{a h_\lambda \otimes b g_\nu}
		\]
		$\sigma$ existiert und setzt $\pi$ fort.
		Also gilt für $x \in J \odot B \subset J \tensormax B$
		\begin{align}
			\norm*{x}_{\max, J \odot B} &= \sup \set[\big]{\norm*{\pi(x)} \given \pi \text{ nicht deg. Darst. von } J \odot B} \\
			&\le \sup \set[\big]{\norm*{\sigma(x)} \given \sigma \text{ nicht deg. Darst. von } A \odot B} \\
			&\le \norm*{\rho(x)} \\
			&\le \norm*{x}_{\max, J \odot B} 
		\end{align} 
		Damit ist $\rho$ isometrisch und $J \tensormax B = J \tensormin B$ für $B$ beliebig.
		Also ist $J$ nuklear.
		$\sfrac{A}{J}$ zeigen wir später.\todo{referenzieren bzw. hier einfügen}
		
		Es seien nun $J$ und $\sfrac{A}{J}$ nuklear. 
		Wir müssen zeigen, dass auch $A$ nuklear ist.
		Betrachte dazu das folgende kommutative Diagramm
		\[
			\begin{tikzcd}
				0 \rar & B \tensormax J \rar \dar["\varphi_J","\cong"'] & B \tensormax A \rar \dar["\varphi_A"] & B \tensormax \sfrac{A}{J} \dar["\varphi_{\sfrac{A}{J}}","\cong"'] \rar & 0 \\
				0 \rar & B \tensormin J \rar & B \tensormin A \rar & B \tensormin \sfrac{A}{J} \rar & 0
			\end{tikzcd}
		\]
		Die obere Zeile ist exakt nach \autoref{prop:24}.
		Wir wollen zeigen, dass auch die untere Zeile exakt ist.
		Wir haben $B \odot A \cap B \tensormin J = B \odot J$ und damit folgt
		\[
			B \odot \faktor{A}{J} = \faktor{B \odot A}{B \odot J} = \faktor{B \odot A}{B \odot A \cap B \tensormin J} \,\,\stackrel[\text{dicht}]{}{\subset}\,\, \faktor{B \tensormin A}{B \tensormin J}
		\]
		Damit induziert $\norm*{\cdot}_{\min}$ eine $C^*$-Norm $\gamma$ auf $B \odot \sfrac{A}{J}$ mit $B \mathop{\otimes_\gamma} \sfrac{A}{J} = \sfrac{B \tensormin A}{B \tensormin J}$.
		Da $\sfrac{A}{J}$ nuklear ist, muss also $\gamma(\cdot) = \norm*{\cdot}_{\min, B \odot \sfrac{A}{J}}$ sein und damit $B \otimes_\gamma \sfrac{A}{J} = B \tensormin \sfrac{A}{J}$.
		Damit ist auch die untere Zeile exakt.
		Nach dem 5er-Lemma ist $\varphi_A$ dann auch ein Isomorphismus.
		\qedhere
	\end{enumerate}
\end{beweis}

\begin{beispiel}
	\begin{enumerate}[(i)]
		\item Ist $A$ endlichdimensional, so ist $A$ auch nuklear.
		\item $A=C_0(X)$ ist nuklear nach \autoref{satz:121}.
		\item Ist $A$ AF\todo{was sollt das sein?}\footnote{z.B. $M_{2^\infty} = M_2 \otimes M_2 \otimes \ldots$} ist nach (i) und \autoref{prop:22}(ii) $A$ nuklear.
		\item Ist $A$ AH, so ist $A$ nuklear. Zum Beispiel
		\[
			A_\theta = C^* \enbrace*{u,v \mid u,v \text{ unitär}, uv=e^{2 \pi i \theta}vu}
		\]
		\item Die Teoplitzalgebra ist nuklear: Nach \autoref{prop:22}(iii) und (iii)
		\[
			\begin{tikzcd}
				0 \rar & \mathcal{K} \rar & \mathcal{T} \rar & C(S^1) \rar & 0
			\end{tikzcd}
		\]
		\item {\large¿} $C_r^*(G)$ (reduzierte Gruppen-$C^*$-Algebra) {\large?} \quad $\leadsto$ \quad {\large¡} Amenabilität {\large!}
		\item {\large¿} $A \rtimes_{\alpha(r)} G$ {\large?}
	\end{enumerate}
\end{beispiel}
\begin{beweis}
	\begin{enumerate}[(i)]
		\item Wegen \autoref{prop:22} (i) können wir $A=M_r$ annehmen.
		Sei $B$ eine $C^*$-Algebra.
		Dann ist $A \odot B$ bereits vollständig bezüglich $\norm*{\cdot}_{\min}$: 
		Sei dazu $(d_\lambda) \subset A \odot B$ ein Cauchy-Netz, wo $d_\lambda = \sum_{i,j=1}^r e_{ij} \otimes d_{\lambda,ij}$ mit $d_{\lambda,ij} \in B$ für $\lambda, i,j$.
		Dann ist
		\[
			e_{ij} \otimes d_{\lambda,ij} = \enbrace*{e_{ii} \otimes \ind_{B^\simm}} d_\lambda \enbrace*{e_{jj} \otimes \ind_{B^\simm}}
		\]
		ebenfalls Cauchy.\footnote{$\varepsilon \leadsto \lambda_0$, für $\lambda_1, \lambda_2 \ge \lambda_0$  $\norm*{d_{\lambda_1,ij} - d_{\lambda_2,ij}} = \norm*{e_{ij} \otimes \enbrace*{d_{\lambda_1,ij}- d_{\lambda_2, ij}}} = \norm*{e_{ij} \otimes d_{\lambda_1,ij}- e_{ij} \otimes d_{\lambda_2, ij}}_{\min} < \varepsilon$}
		Aber $\norm*{e_{ij} \otimes b}_{\min} = \norm*{e_{ij}}_A \norm*{b}_B = \norm*{b}_B$.
		Also ist $(d_{\lambda,ij})_{\Lambda} \subset B$ Cauchy und konvergiert in $B$ für jedes $i,j$.
		Damit konvergiert auch $(d_\lambda)$ und somit ist $\enbrace*{A \odot B, \norm*{\cdot}_{\min}}$ eine $C^*$-Algebra.
		Daraus folgt $\norm*{\cdot}_{\min}=\norm{\cdot}_{\max}$ mit der Eindeutigkeit einer \emph{vollständigen} $C^*$-Norm.\qedhere
	\end{enumerate}
\end{beweis}

\begin{proposition}[label=prop:24,{name=[{das maximale Tensorprodukt ist exakt}]}]
	$A \tensormax -$ ist ein exakter Funktor für jede $C^*$-Algebra $A$
\end{proposition}
\begin{beweis}
	Für Funktorialität benutze \autoref{prop:125}.
	$\pi \colon B \to C$ induziert also $\id_A \tensormax \pi \colon A \tensormax B \to A \tensormax C$.
	Sei nun \(
		\begin{tikzcd}[cramped,sep=small]
			0 \rar & J \rar["\iota"] & B \rar["\pi"] & \sfrac{B}{J} \rar & 0
		\end{tikzcd}
	\) exakt.
	Zu zeigen ist, dass folgende Seqenz exakt ist
	\[
		\begin{tikzcd}[column sep=3.8em]
			0 \rar & A \tensormax J \rar["\id_A \tensormax \iota"] & A \tensormax B \rar["\id_A \tensormax \pi"] & A \tensormax \sfrac{B}{J} \rar & 0
		\end{tikzcd}
	\]
	$\id_A \tensormax \pi$ hat dichtes Bild und vollständiges Bild und ist damit surjektiv, womit die Exaktheit an der dritten Stelle gezeigt ist.
	
	Für die Injektivität von $\id_A \tensormax \iota$ sei $f \in S(A \tensormax J) \StackText{\ref{satz:119}}{\approx} S(A \odot J)$ und $(u_\lambda)$ eine approximative Eins für $J$.
	Dann definiert $a \otimes b \mapsto \lim_\lambda f \enbrace*{a \otimes u_\lambda b u_\lambda}$ eine Fortsetzung von $f$ zu $\tilde{f} \in S(A \odot B)\approx S(A \tensormax B)$ (warum?).
	Nach \autoref{satz:119} ist nun $\norm*{\cdot}_{A \tensormax J} \big|_{A \odot J} = \norm*{\cdot}_{A \tensormax B}\big|_{A \odot J}$, also ist ${\id_{A}} \tensormax \iota \big|_{A \odot J}$ isometrisch.
	Wegen Stetigkeit ist dann auch ${\id_A} \tensormax \iota$ isometrisch.
	
	Zur Exaktheit an der zweiten Stelle: 
	Wir haben $A \tensormax J \lhd A \tensormax B$ (warum?) und außerdem
	\[
		A \tensormax J \subset \ker \enbrace*{{\id_A} \tensormax \pi}
	\]
	Damit erhalten wir eine induzierte Abbildung $\sigma\colon\sfrac{A \tensormax B}{A \tensormax J} \twoheadrightarrow \sfrac{A \tensormax B}{\ker \enbrace*{{\id_A} \tensormax \pi}} = A \tensormax \sfrac{B}{J}$.
	Weiter haben wir $A \odot B \cap A \tensormax J = A \odot J$.
	Damit gilt
	\[
		A \odot \faktor{B}{J} = \faktor{A \odot B}{A \odot J} = \faktor{A \odot B}{A \odot B \cap A \tensormax J} \,\,\stackrel[\text{dicht}]{}{\subset}\, \, \faktor{A \tensormax B}{A \tensormax J}
	\]
	Also erhalten wir $\rho \colon A \tensormax \sfrac{B}{J} \twoheadrightarrow \sfrac{A \tensormax B}{A \tensormax J}$ mit $\sigma \colon \circ \rho \big|_{A \odot \sfrac{B}{J}} = \id_{A \odot \sfrac{B}{J}}$.
	$\sigma \circ \rho$ ist also isometrisch auf $A \odot \sfrac{B}{J}$ und wegen Stetigkeit damit überall.
	Da $\rho$ surjektiv ist, muss $\sigma$ insbesondere injektiv sein.
	Also gilt $A \tensormax J = \ker \enbrace*{{\id_A} \otimes \pi}$ und damit Exaktheit.
\end{beweis}

\begin{definition}[{name=[exakte C^*-Algebra]}]
	Eine $C^*$-Algebra $A$ heißt \bet{exakt}\index{exakte $C^*$-Algebra}, falls $A \tensormin -$ ein exakter Funktor ist, das heißt falls \(
		\begin{tikzcd}[cramped,sep=small]
			0 \rar & J \rar["\iota"] & B \rar["\pi"] & \sfrac{B}{J} \rar & 0
		\end{tikzcd}
	\)
	exakt ist, so ist auch die folgende Sequenz exakt
	\[
		\begin{tikzcd}
			0 \rar & A \tensormin J \rar["{\id_A} \otimes \iota"] & A \tensormin B \rar["{\id_A} \otimes \pi"] & A \tensormin \sfrac{B}{J} \rar & 0
		\end{tikzcd}
	\]
\end{definition}

\begin{bemerkung}[{name=[{nukleare $C^*$-Algebren sind exakt}]}]
	Nach \autoref{prop:24} ist jede nukleare $C^*$-Algebra exakt.
\end{bemerkung}

\begin{satz}[{name={Kirchberg, Wassermann}}]
	Eine $C^*$-Unteralgebra einer exakten $C^*$-Algebra ist exakt.
	Insbesondere sind $C^*$-Unteralgebren von nuklearen $C^*$-Algebren exakt.
\end{satz}
\begin{beweis}[Idee]
	Es seien $A \subset C$ $C^*$-Algebren und $C$ exakt.
	Sei \(
		\begin{tikzcd}[cramped,sep=small]
			0 \rar & J \rar["\iota"] & B \rar["\pi"] & \sfrac{B}{J} \rar & 0
		\end{tikzcd}
	\) eine beliebige exakte Sequenz.
	Es genügt zu zeigen: 
	\[
		A \tensormin J = \ker \enbrace*{{\id_A} \tensormin \pi}
	\]
	Dabei ist die Inklusion \enquote{$\subseteq$} klar.
	Die andere zeigt man in mehereren Schritten:
	\begin{enumerate}[1.]
		\item Es gilt $C \tensormin J \cap A \tensormin B = A \tensormin J$.
		
		Dabei ist \enquote{$\supseteq$} klar. 
		Die andere Inklusion zeigt man mit Hilfe einer approximativen Eins.
		\item Sei $\varphi \in S(A)$. 
		Man benutzt die \emph{slice map}
		\mapdef{\varphi \otimes {\id_B} \colon A \odot B}{\mathbb{C} \odot B = B}{a \otimes b}{\varphi(a) \cdot b}{}
		$\varphi \otimes {\id_B}$ ist stetig bezüglich $\norm*{\cdot}_{{\min}, A \odot B}$ und wir erhalten eine stetige Fortsetzung $\varphi \tensormin  {\id_B} \colon A \tensormin B \to B$. (Übung)\todo{\TeX{}en?}
		\item Betrachte 
		\[
			K_J \coloneqq \set[\big]{x \in A \tensormin B \given (\varphi \otimes {\id_B})(x^* x) \in J \, \forall \varphi \in S(A)}
		\]
		und zeige $K_J = \ker \enbrace*{{\id_A} \tensormin \pi_n}$ (benutze \autoref{satz:120}).
		\item Es gilt $K_J \subset A \tensormin J$:
		Sei $x_0 \in K_J$.
		Falls $\psi \in S(C)$, so ist $\varphi \coloneqq \psi|_A$ bis auf Normierung ein Zustand auf $A$ und es gilt $(\psi \otimes {\id_B})(x_0^*x_0) = (\varphi \otimes {\id_B})(x_0^* x_0) \in J$.
		Also gilt
		\[
			x_0 \in \set*{x \in C \tensormin B \given (\rho \otimes {\id_B})(x^*x) \in J \, \forall \rho \in S(C)} \StackText{3.}{=} \ker \enbrace*{{\id_C} \tensormin \pi_n} = C \tensormin J
		\]
		Damit ist $x_0 \in \enbrace*{A \tensormin B} \cap \enbrace*{C \tensormin J} \StackText{1.}{=} A \tensormin J$.\qedhere
	\end{enumerate}
\end{beweis}
% section 2 (end)
\newpage

\section{Vollständig positive Abbildungen} % (fold)
\label{sec:3}

\begin{erinnerungA}[{name=[{induzierte Abbildung auf Matrizen}]}]
	Eine Abbildung $\varphi\colon A \to B$ induziert eine Abbildung $\varphi^\ssbrace{n} \colon M_n(A) \to M_n(B)$ durch $(a_{ij})_{ij} \mapsto \enbrace*{\varphi(a_{ij})}_{ij}$.
	Dabei müssen $A$ und $B$ keine $C^*$-Algebren sein.
	Sind $A$ und $B$ Vektorräume, so ist $\varphi^\ssbrace{n} = \varphi \otimes \id_{M_n}$ und aus Linearität von $\varphi$ folgt Linearität von $\varphi^\ssbrace{n}$.
	Ist $\varphi$ ein $^*$-Homomorphismus, so ist auch $\varphi^\ssbrace{n}$ ein $^*$-Homomorphismus.
	Wir wir in \autoref{bsp:35} sehen werden, muss $\varphi^\ssbrace{n}$ aber nicht positiv sein, wenn $\varphi$ positiv ist!
	Dies führt uns zu folgender Definition:
\end{erinnerungA}

\begin{definition}[{name=[positiv, vollständig positiv und beschränkt]},label=def:32]
	Seien $A,B$ $C^*$-Algebren und $\varphi \colon A \to B$ linear
	\begin{enumerate}[(i),itemsep=1pt]
		\item $\varphi$ heißt \bet{positiv}, falls $\varphi(A_+) \subset B_+$.\index{positive Abbildung}
		\item $\varphi$ heißt $\mathbold n$\bet{-positiv}, falls $\varphi^\ssbrace{n}$ positiv ist.\index{n-positive Abbildung@$n$-positive Abbildung}
		\item $\varphi$ heißt \bet{vollständig positiv}, falls $\varphi$ $n$-positiv ist für alle $n \in \mathbb{N}$.\index{vollständig positive Abbildung}
		\item $\varphi$ heißt \bet{vollständig beschränkt}, falls 
		\[
			\norm*{\varphi}_\vb \coloneqq \sup_{n \in \mathbb{N}} \norm*{\varphi^\ssbrace{n}} < \infty
		\]
 	\end{enumerate}
\end{definition}

\begin{definition}[{name=[Operatorsystem]}]
	Sei $A$ eine unitale $C^*$-Algebra. 
	Ein unitaler selbstadjungierter Untervektorraum von $X$ heißt \Index{Operatorsystem}. 
\end{definition}

\begin{bemerkung}[{name=[grundlegende Eigenschaften von Operatorsystemen]}]
	\leavevmode
	\begin{enumerate}[(i)]
		\item Sei $X$ ein Operatorsystem.
		Dann gilt $X = \Span X_+$:
		
		Ist $x \in X$, so sind $x+x^*$ sowie $ i(x-x^*)$ in $X_\sa$ enthalten und es folgt
		\[
			x = \sfrac{1}{2}\enbrace[\big]{x + x^* - i \cdot i (x-x^*)} \implies x \in \Span X_\sa
		\]
		Ist $y \in X_\sa$, so sind $\norm*{y} \cdot \ind -y, \norm*{y} \cdot \ind \in X_+$ und es folgt
		\[
			y = \norm*{y} \cdot \ind - \enbrace[\big]{\norm*{y} \cdot \ind -y} \implies X_\sa = \Span X_+
		\]
		\item \autoref{def:32} ist auch sinnvoll für $\varphi \colon X \to Y$ linear, wobei $X$ und $Y$ Operatorsysteme sind.
	\end{enumerate}
\end{bemerkung}

\begin{beispiel}[label=bsp:35,{name=[{n-positive und vollständig positive Abbildungen}]}]
	\leavevmode
	\begin{enumerate}[(i)]
		\item Ist $\varphi \colon A \to B$ ein $^*$-Homomorphismus, so ist $\varphi$ vollständig positiv.
		Die Positivität folgt mit
		\[
			\varphi(x^*x) = \varphi(x)^* \varphi(x) \ge 0
		\]
		Da $\varphi^\ssbrace{n}$ auch ein $^*$-Homorphismus ist, ist auch $\varphi^\ssbrace{n}$ positiv für alle $n \in \mathbb{N}$.
		\item Sei $\varphi \colon A \to B$ $n$-positiv und $v \in B$ fest gewählt.
		Dann ist die \Index{Kompression}
		\mapdef{\varphi_v \colon A}{B}{a}{v^* \varphi(a) v}{}
		ebenfalls $n$-positiv:
		Es ist $(\varphi_v)^\ssbrace{n} = \enbrace*{\varphi^\ssbrace{n}}_{v \otimes \ind_{M_n}}$ und nach \cite[Satz~3.7 (v)]{OpAlg} gilt
		\[
			(\varphi_v)^\ssbrace{n}(x^* x) = \Underbracket{\enbrace*{v^* \otimes \ind_{M_n}} \Underbracket{\varphi^\ssbrace{n}(x^*x)}{\ge 0} \enbrace*{v \otimes \ind_{M_n}}}{\ge 0}
		\]
		Insbesondere gilt: Ist $\varphi$ vollständig positiv, so ist $\varphi_v$ vollständig positiv. Ebenso wenn $\varphi$ ein $^*$-Homomorphismus ist.
		Wir werden weiter sehen: \emph{Jede vollständig positive Abbildung ist von dieser Form!}\todo{referenzieren!}
		\item $\varphi \colon \mathbb{C}^k \to B$ ist positiv genau dann, wenn $b_i \coloneqq \varphi(e_i) \ge 0$ für $i=1,\ldots ,k$.
		
		$\varphi \colon \mathbb{C}^k \to C(X)$ ist unital und positiv genau dann, wenn $(b_i)_{i=1,\ldots ,k}$ Zerlegung der Eins ist.
		\item Die Involution $\tau \colon M_2 \to M_2$, $\tau(a)=a^*$ ist positiv, aber nicht $2$-positiv:
		Die Positivität folgt sofort mit $\tau(a^*a)= a^*a \ge 0$. Betrachtet man die Matrix von Matrizen\todo{RevChap 3}
		\[
			x \coloneqq \begin{pmatrix}
				1 & 0 & 0 & 1 \\
				0 & 0 & 0 & 0 \\
				0 & 0 & 0 & 0 \\
				1 & 0 & 0 & 1
			\end{pmatrix} \in M_2(M_2)
		\]
		so gilt $x \ge 0$ (warum?), aber 
		\[
			\tau^\ssbrace{2}(x) = \begin{pmatrix}
				1 & 0 & 0 & 0 \\
				0 & 0 & 1 & 0 \\
				0 & 1 & 0 & 0 \\
				0 & 0 & 0 & 1
			\end{pmatrix} \not\ge 0
		\]
		(warum?)
	\end{enumerate}
\end{beispiel}

\begin{proposition}[label=prop:36]
	Seien $A,B$ $C^*$-Algebren und $\varphi \colon A \to B$ positiv.
	Dann ist $\varphi$ beschränkt.
\end{proposition}
\begin{beweis}
	Falls $\varphi$ unbeschränkt auf $A$ ist, so auch auf $A_+$.
	Folglich existiert $(a_n)_{\mathbb{N}} \subset A_+^1$ mit $\norm*{\varphi(a_n)} \ge n^3$ für $n \in \mathbb{N}$.
	Dann ist $a \coloneqq \sum_{n=0}^{\infty} \sfrac{1}{n^2} a_n \in A$ und $\sfrac{1}{n^2} a_n \le a$ für $n \in \mathbb{N}$.
	Damit haben wir für alle $n \in \mathbb{N}$
	\[
		n \le \sfrac{1}{n^2} \cdot \norm*{\varphi(a_n)} = \norm*{\varphi \enbrace*{\sfrac{1}{n^2} \cdot a_n}} \le \norm*{\varphi(a)} 
	\]
	Dies ist ein Widerspruch.
\end{beweis}

\begin{bemerkung}[label=bem:37]
	\leavevmode
	\begin{enumerate}[(i)]
		\item Sei $X$ ein Operatorsystem und $\varphi \colon X \to B$ positiv.
		Dann gilt $\norm*{\varphi} \le 2 \cdot \norm*{\varphi(\ind)}$ (warum?)
		\item In (i) kann Gleichheit auftreten: 
		Betrachte dazu $X \coloneqq \Span \set*{1,z,\overline{z}} \subset C(\mathbb{T})$ und $\varphi \colon X \to M_2$ mit
		\[
			\varphi \enbrace*{a \cdot 1 + b \cdot z + c \cdot \overline{z}} = \begin{pmatrix}
				a & 2 b \\
				2 c & a
			\end{pmatrix}
		\]
		$\varphi$ ist positiv (warum?).
		Es gilt nach (i) $2 \cdot \norm*{\varphi(1)} = 2 = \norm*{\varphi(z)} \le \norm*{\varphi} \le 2 \cdot \norm*{\varphi(1)}$.
	\end{enumerate}
\end{bemerkung}

\begin{lemma}[label=lem:38]
	Sei $A$ eine $C^*$-Algebra und $a,h \in A$ ($h \in A_\sa$).
	Dann gilt
	\begin{enumerate}[(i)]
		\item $\begin{psmallmatrix}
			\ind_{A^\simm} & a \\
			a^* & \ind_{A^\simm}
		\end{psmallmatrix} \ge 0 \iff \norm*{a} \le 1$
		\item $\begin{psmallmatrix}
			h & a \\
			a^* & h
		\end{psmallmatrix} \ge 0 \implies a^* \le \norm*{h} \cdot h$, also $\norm*{a} \le \norm*{h}$\marginnote{Die Umkehrung gilt nicht!}
		\item $\begin{psmallmatrix}
			\ind_{A^\simm} & a \\
			a^* & h
		\end{psmallmatrix} \ge 0 \iff a^*a \le h$. 
	\end{enumerate}
\end{lemma}
\begin{beweis}
	Ohne Einschränkungen können wir $A \subset \mathcal{B}(\mathcal{H})$ und $M_2(A) \subset M_2(\mathcal{B}(\mathcal{H})) = \mathcal{B}(\mathcal{H} \oplus \mathcal{H})$ annehmen.
	\begin{enumerate}[(i)]
		\item Es gilt
		\begin{align}
			\skal*{\binom{\xi}{\eta}}{ \begin{pmatrix}
				\ind_{A^\simm} & a \\
				a^* & \ind_{A^\simm}
			\end{pmatrix} \binom{\xi}{\eta}} &= \skal*{\xi}{\xi} + \skal*{\xi}{a \eta} + \skal*{\eta}{a^* \xi} + \skal*{\eta}{\eta} \label{eq:38:1} \tag{$*$}\\
			&\ge \skal*{\xi}{\xi} - 2 \norm*{a} \cdot \norm*{\xi} \cdot \norm*{\eta} + \skal*{\eta}{\eta} \label{eq:38:2} \tag{$**$}
		\end{align}
		Falls $\norm*{a} \le 1$ ist also \eqref{eq:38:2} $\ge 0$ und somit $\begin{psmallmatrix}
			1 & a \\ a^* & 1
		\end{psmallmatrix} \ge 0$.
		Falls $\norm*{a} > 1$, so existieren $\xi,\eta \in \mathcal{H}$ mit $\norm*{\xi} = \norm*{\eta}=1$ und $\skal*{\xi}{a \eta} < -1$
		Damit folgt \eqref{eq:38:2} $<0$ und somit $\begin{psmallmatrix}
			1 & a \\ a^* & 1
		\end{psmallmatrix} \not\ge 0$.
		\item Angenommen $h \neq 0$. Für $\eta \in \mathcal{H}$ setzen wir $\xi \coloneqq - \sfrac{1}{\norm*{h}} \cdot a \cdot \eta$.
		Dann gilt
		\[
			\skal*{\binom{\xi}{\eta}}{\begin{pmatrix}
				h & a \\ a^* & h
			\end{pmatrix} \binom{\xi}{\eta}} = \ldots \le \sfrac{1}{\norm*{h}} \skal*{\eta}{a^* a \eta} - \sfrac{2}{\norm*{h}} \skal*{\eta}{a^* a \eta} + \skal*{\eta}{h \eta}
		\]
		und daraus folgt $\skal*{\eta}{ \enbrace*{h - \sfrac{1}{\norm*{h}} a^*a}\eta} \ge 0$ für $\eta \in \mathcal{H}$, also $h \ge \sfrac{1}{\norm*{h}} a^* a$.\marginnote{für $h=0$ den Bruch jeweils weglassen}
	\end{enumerate}
\end{beweis}

\begin{proposition}[label=prop:39]
	Sei $X$ ein Operatorsystem, $B$ eine unitale $C^*$-Algebra und $\varphi \colon X \to B$ unital und $2$-positiv.
	Dann gilt $\norm*{\varphi}=1$.
\end{proposition}
\begin{beweis}
	Sei $a \in X$ mit $\norm*{a} \le 1$. Mit \autoref{lem:38} (i) folgt 
	\[
		\begin{pmatrix}
			\ind_{X} & a \\ a^* & \ind_X
		\end{pmatrix} \ge 0 \implies \begin{pmatrix}
			\ind_B & \varphi(a) \\ \varphi(a)^* & \ind_B
		\end{pmatrix} \ge 0
	\]
	woraus mit \autoref{lem:38} (i) wieder $\norm*{\varphi(a)} \le 1$ folgt.
\end{beweis}

\begin{proposition}[{name={Cauchy-Schwarz}},label=prop:310]
	Es seien $A$ und $B$ unitale $C^*$-Algebren und $\varphi \colon A \to B$ unital und $2$-positiv.
	Dann gilt für $a \in A$
	\[
		\varphi(a)^* \varphi(a) \le \varphi(a^* a)
	\]
\end{proposition}
\begin{beweis}
	Es gilt
	\[
		\begin{pmatrix}
			\ind_A & a \\ a^* & a^* a
		\end{pmatrix} 
		=
		\begin{pmatrix}
			\ind_A & 0  \\ a^* & 0
		\end{pmatrix} \cdot 
		\begin{pmatrix}
			\ind_A & a \\ 0 & 0
		\end{pmatrix} \ge 0
	\]
	Da $\varphi$ $2$-positiv ist, folgt
	\[
		\begin{pmatrix}
			\ind_B & \varphi(a) \\ \varphi(a)^* & \varphi(a^* a) 
		\end{pmatrix} \ge 0
	\]
	woraus wir mit \autoref{lem:38} (ii) $\varphi(a)^* \varphi(a) \le \varphi(a^* a)$ erhalten.
\end{beweis}

\begin{definitionP}
	Sei $A$ eine $C^*$-Algebra.
	Dann ist
	\[
		\mathcal{F}(A) \coloneqq \set[\big]{a \in A \given \exists e \in A : \norm*{e}\le 1, a e = e a = a}
	\]
	dicht in $A$.
	Weiter ist $\mathcal{F}(A)_+ \subset A_+$ dicht und $\mathcal{F}(A)_\sa \subset A_\sa$ dicht.
\end{definitionP}
\begin{beweis}
	Sei $(u_\lambda)_{\Lambda}$ eine approximative Eins.
	Für $\varepsilon>0$ definiere $f_\varepsilon $ und $g_\varepsilon$ durch
	\missingfigure{Hier fehlen noch zwei Bilder von Funktionsgraphen}
	Es gilt $\lim_{\lambda, \varepsilon} f_\varepsilon(u_\lambda) a f_\varepsilon(u_\lambda) = a$ für alle $a \in A$ und $f_\varepsilon(u_\lambda) a f_\varepsilon(u_\lambda) \in \mathcal{F}(A)_{(+, \sa)}$, da $g_\varepsilon(u_\lambda) f_\varepsilon(u_\lambda) = f_\varepsilon(u_\lambda)$.
\end{beweis}

\begin{definition}[{name=[{Pedersen-Ideal}]}]
	\[
		\Ped(A) \coloneqq A \, \mathcal{F}(A)\, A \stackrel[\text{dicht}]{}{\subset} A
	\]
	ist das (eindeutig bestimmte) minimale dichte Ideal von $A$.\marginnote{$A$ unital $\Rightarrow$ $\Ped(A)=A$}
	$\Ped(A)$ heißt das \Index{Pedersen-Ideal}. 
\end{definition}

\begin{proposition}[label=prop:313]
	Seien $A,B$ beides $C^*$-Algebren, $\varphi \colon A \to B$ vollständig positiv und $\mu \ge \norm*{\varphi}$.
	Sei $\varphi^+ \colon A^+ \to B^+$ (bzw. $B^\simm$) die durch $\varphi^*(\ind_{A^+}) \coloneqq \mu \cdot \ind_{B^+}$ definierte lineare Fortsetzung.
	Dann ist $\varphi^+$ vollständig positiv und es gilt $\norm*{\varphi^+}= \mu$.
\end{proposition}
\begin{beweis}
	Wir müssen zeigen, dass $(\varphi^+)^\ssbrace{n}$ positiv ist: 
	Sei $\pi \colon A^+ \to \mathbb{C}$ der kanonische $^*$"=Homomorphismus.
	Dann ist auch $\pi^\ssbrace{n} \colon M_n(A^+) \to M_n$ ein $^*$-Homomorphismus, also insbesondere positiv.
	Sei $0 \le a + x  \otimes \ind_{A^+} \in M_n(A^+) = M_n \otimes A^+$ mit $a=a^* \in M_n(A)$, $x \ge 0 \in M_n$.\footnote{$\sum y_i \otimes a_i= \sum y_i \otimes (b_i + \lambda \cdot \ind_{A^+})$}
	Zu zeigen: 
	\[
		0 \le (\varphi^+)^\ssbrace{n} \enbrace*{a + x \otimes \ind_{A^+}} = \varphi^\ssbrace{n}(a) + \mu x \otimes \ind_{B^+}
	\]
	Es genügt dies für $a \in \mathcal{F}(M_n(A))_\sa$ zu zeigen, da $\mathcal{F}(M_n(A))_\sa \subset M_n(A)_\sa$ dicht, $\varphi^\ssbrace{n}$ beschränkt und $M_n(B^+)_+ \subset M_n(B_+)$ abgeschlossen ist (Warum genau?).
	Wir dürfen sogar annehmen, dass $e \in A_+$ existiert mit $\norm*{e}\le 1$ und $\overline{e}a = a \overline{e} =a$, wo $\overline{e}= \diag(e,\ldots ,e)=\ind_n \otimes e$.
	Dann ist
	\begin{align}
		(\varphi^+)^\ssbrace{n} (a + x \otimes \ind_{A^+}) &= (\varphi^+)^\ssbrace{n} \enbrace[\big]{\overline{e} \enbrace*{a + x \otimes \ind_{A^+}}\overline{e} + (1- \overline{e}^2) \enbrace*{ a + x \otimes \ind_{A^+}}} \\
		&= (\varphi^+)^\ssbrace{n} \enbrace[\Big]{\Underbracket{\overline{e} \enbrace*{a + x \otimes \ind_{A^+}} \overline{e}}{\ge 0 \in M_n(A)}} + (\varphi^+)^\ssbrace{n} \enbrace*{\Underbracket{x \otimes (1 - e^2)}{\ge 0 }} \\
		&\ge 0
	\end{align}
\end{beweis}

\begin{proposition}[label=prop:314]
	Seien $A,B$ $C^*$-Algebren und $\varphi \colon A \to B$ vollständig positiv.
	Dann ist $\varphi$ vollständig beschränkt und es gilt für eine approximative Eins $(u_\lambda)_\Lambda \subset A$
	\[
		\norm*{\varphi}_\vb = \norm*{\varphi} = \sup_\lambda \norm*{\varphi(u_\lambda)}
	\]
\end{proposition}
\begin{beweis}
	Es gilt offenbar $\norm*{\varphi(u_\lambda)} \le \norm*{\varphi} \le \norm*{\varphi}_\vb$.
	Zu zeigen bleibt $\norm*{\varphi}_\vb \le \sup_\lambda \norm*{\varphi(u_\lambda)}$.
	Sei also $a \in M_n(A)$ mit $\norm*{a} \le 1$ und $\overline{u}_\lambda = \diag (u_\lambda,\ldots ,u_\lambda)$.
	Dann ist
	\[
		0 \StackText{\ref{lem:38}(i)}{\le} \begin{pmatrix}
			\ind_{M_n(A)^\simm} & a \\ a^* & \ind_{M_n(A)^\simm}
		\end{pmatrix}
		\implies
		0 \le \begin{pmatrix}
			\overline{u}_\lambda^{\sfrac{1}{2}} & 0 \\ 0 & \overline{u}_\lambda^{\sfrac{1}{2}}
		\end{pmatrix}
		\begin{pmatrix}
			\ind & a \\ a^* & \ind
		\end{pmatrix}
		\begin{pmatrix}
			\overline{u}_\lambda^{\sfrac{1}{2}} & 0 \\ 0 & \overline{u}_\lambda^{\sfrac{1}{2}}
		\end{pmatrix}
	\]
	also gilt da $\varphi$ vollständig positiv ist
	\[
		0 \le \begin{pmatrix}
			\varphi^\ssbrace{n}(\overline{u}_\lambda) & \varphi^\ssbrace{n} \enbrace*{\overline{u}_\lambda^{\sfrac{1}{2}} a \overline{u}_\lambda^{\sfrac{1}{2}}} \\
			\varphi^\ssbrace{n}\enbrace*{\overline{u}_\lambda^{\sfrac{1}{2}} a^* \overline{u}_\lambda^{\sfrac{1}{2}}} & \varphi^\ssbrace{n}(\overline{u}_\lambda)
		\end{pmatrix}
	\]
	Mit \autoref{lem:38} (iii) folgt nun 
	\[
		\Underbracket{\norm*{\varphi^\ssbrace{n} \enbrace*{\overline{u}_\lambda^{\sfrac{1}{2}} a \overline{u}_\lambda^{\sfrac{1}{2}}}}}{\grenzw{} \norm*{\varphi^\ssbrace{n} (a)}} \le \norm*{\varphi^\ssbrace{n} (\overline{u}_\lambda)} = \norm*{\varphi(u_\lambda)} \grenzw{\lambda} \sup_\lambda \norm*{\varphi(u_\lambda)} \qedhere
	\]
\end{beweis}

\begin{bemerkung}[label=bem:315]
	Die Aussage gilt entsprechend, falls $A$ ein Operatorsystem ist.
\end{bemerkung}

\begin{proposition}[label=prop:316]
	Sei $A$ eine $C^*$-Algebra und $\varphi \colon A \to \mathbb{C}$ positiv.
	Dann ist $\varphi$ vollständig positiv und es gilt.
	Ebenso für $\varphi \colon X \to \mathbb{C}$ positiv, falls $X$ ein Operatorsystem ist.
\end{proposition}
\begin{beweis}
	Betrachte $0 \le (a_{ij})_{ij} \in M_n(A)$ und $\xi= (\xi_1, \ldots ,\xi_n)^T \in \mathbb{C}^n$.
	Dann gilt
	\begin{align}
		\skal[\big]{\xi}{\varphi^\ssbrace{n} \enbrace*{(a_{ij})_{ij}}} = \sum_{i,j=1}^n \varphi(a_{ij}) \overline{\xi}_i \xi_j &= \varphi \enbrace*{\sum_{i,j=1}^n \overline{\xi}_i \xi_j a_{ij}} \ge 0 \\
		&= 
	\end{align}
	Dann ist 
	\[
		M_n(A) \ni \begin{pmatrix}
					\overline{\xi}_1 \ind_{A^\simm} & \cdots & \overline{\xi}_n \ind_{A^\simm} \\
					& &  \\
					& 0 & 
				\end{pmatrix}\enbrace[\big]{a_{ij}}_{ij} 
				\begin{pmatrix}
					\xi_1 \ind_{A^\simm} & & \\
					\vdots & & 0 \\
					\xi_n \ind_{A^\simm}  & & 
				\end{pmatrix} \ge 0
	\]
\end{beweis}

\begin{korollar}[label=kor:317]
	Sei $\varphi \colon A \to C_0(\Omega)$ positiv, wobei $\Omega$ lokal kompakt und Hausdorff.
	Dann ist $\varphi$ vollständig positiv.
	Ebenso für $\varphi \colon X \to C_0(\Omega)$, wobei $X$ ein Operatorsystem ist.
\end{korollar}
\begin{beweis}
	Es gilt
	\[
		\varphi^\ssbrace{n} \enbrace[\big]{(a_{ij})_{ij}} \ge 0 \iff \varphi^\ssbrace{n} \enbrace[\big]{(a_{ij})_{ij}}(t) \ge 0 , \enspace t \in \Omega
	\]
	Aber $\varphi^\ssbrace{n}(\cdot )(t) = \enbrace*{\varphi(\cdot )(t)}^\ssbrace{n}$ ist vollständig positiv nach \autoref{prop:316}.
\end{beweis}

\begin{proposition}[label=prop:318]
	Sei $\Omega$ ein lokalkompakter Hausdorffraum und $B$ eine $C^*$-Algebra.
	Falls $\varphi \colon C_0(\Omega) \to B$ positiv ist, so ist $\varphi$ vollständig positiv.
\end{proposition}
\begin{beweis}
	Wähle $0 \le a \in M_n(C_0(\Omega))\cong M_n \otimes C_0(\Omega) \cong C_0(\Omega,M_n)$.
	Zu zeigen ist nun, dass $M_n(B) \ni \varphi^\ssbrace{n}(a) \ge 0$.
	Zu $\varepsilon>0$ existieren $0 \le f_1, \ldots ,f_k \in  C_0(\Omega)$ und $0 \le a_1, \ldots ,a_k \in M_n(A)$ mit
	\[
		\norm*{a - \sum\nolimits_{i=1}^{k} a_i \otimes f_i} < \varepsilon
	\]
	(warum?)
	Dann gilt 
	\begin{align}
		\varphi^\ssbrace{n}(a) \ge \varphi^\ssbrace{n} \enbrace*{\sum_{i=1}^k a_i \otimes f_i} - \varepsilon \cdot \norm*{\varphi} \cdot \ind_{M_n \otimes B^\simm} &= \sum_{i=1}^{k} \Underbracket{a_i}{\ge 0} \otimes  \Underbracket{\varphi(f_i)}{\ge 0} - \varepsilon \cdot \norm*{\varphi^\ssbrace{n}} \cdot \ind_{M_n \otimes B^\simm} \\
		&\ge 0 - \varepsilon \cdot \norm*{\varphi^\ssbrace{n}} \cdot \ind
	\end{align}
	Damit ist $\varphi^\ssbrace{n}(a) \ge 0$, da $\varepsilon$ beliebig war und $M_n(B)_+ \subset M_n(B)$ abgeschlossen.
\end{beweis}

\begin{satz}[label=satz:319,name={Stinespring}]
	Sei $A$ eine $C^*$-Algebra, $\varphi \colon A \to \mathcal{B}(\mathcal{H})$ vollständig positiv.
	Dann existieren ein Hilbertraum $\tilde{\mathcal{H}}$, ein $^*$-Homomorphismus $\pi \colon A \to \mathcal{B}(\tilde{\mathcal{H}})$ und ein $v \in \mathcal{B}(\mathcal{H},\tilde{\mathcal{H}})$ mit $\norm*{\varphi} = \norm*{v}^2$ und für $a \in A$\marginnote{$\begin{psmallmatrix}
		0 & 0 \\ v & 0
	\end{psmallmatrix} \in \mathcal{B}(\mathcal{H} \oplus \tilde{\mathcal{H}})$}
	\[
		\varphi(a) = v^* \pi(a) v
	\]
	Das heißt $v$ ist Kompression eines $^*$-Homomorphismus.
	Falls $A,\varphi$ unital sind, so ist $v$ eine Isometrie.
\end{satz}
\begin{beweis}
	GNS-Konstruktion
\end{beweis}
\todo[inline]{Hier fehlt leider noch eine komplette Vorlesung!}
% section 3 (end)




\cleardoubleoddemptypage
\pagenumbering{Alph}
\setcounter{page}{1}
\cleardoubleoddemptypage
\appendix

\section{Anhang} % (fold)
\label{sec:anhang}
%!TEX root = ana_top_geo.tex

\subsection{Ausführlicher Beweis zu \cref{lem:kpt-schnitte}} % (fold)
\label{sub:kpt-schnitte}
Sei $X$ ein Hausdorffraum. Dann ist $X$ genau dann kompakt, wenn gilt: Hat eine Familie $\mathcal{A}$ von abgeschlossenen Teilmengen von $X$ die endliche 
Durchschnittseigenschaft, so gilt 
\[
	\bigcap_{A \in \mathcal{A}} A \not= \emptyset.
\]
\begin{beweis}
	Für die erste Implikation sei $X$ kompakt und $\mathcal{A}$ eine Familie von abgeschlossenen Mengen mit der endlichen Durchschnittseigenschaft.
	Angenommen $\bigcap_{A \in \mathcal{A}} A = \emptyset$.
	Dann gilt
	\[
		X = X \setminus \bigcap_{A \in \mathcal{A}} A = \bigcup_{A \in \mathcal{A}} X \setminus A.
	\]
	Nun ist $\mathcal{U} \coloneqq \set*{X \setminus A \given A \in \mathcal{A}}$ eine offene Überdeckung von $X$ und da $X$ kompakt ist, existiert $\mathcal{A}_0 \subset \mathcal{A}$ endlich, sodass
	\[
		X = \bigcup_{A \in \mathcal{A}_0} X \setminus A = X \setminus \underbrace{\bigcap_{A \in \mathcal{A}_0 } A }_{\neq \emptyset} \quad \light
	\]
	Für die umgekehrte Implikation sei nun $\mathcal{U} = \set{U_i}_{i \in I}$ eine offene Überdeckung von $X$.
	Angenommen für jede endliche Teilmenge $J \subseteq I$ gilt $X \neq \bigcup_{i \in J} U_i$.
	Betrachte nun $\mathcal{A} =  \set{X \setminus U_i}_{i \in I}$. Dann gilt nach Annahme
	\[
		\bigcap_{i \in J} X \setminus U_i = X \setminus \bigcup_{i \in J} U_i \neq \emptyset.
	\]
	Also hat $\mathcal{A}$ die endliche Durchschnittseigenschaft. Nach Vorraussetzung gilt dann
	\[
		\emptyset \not= \bigcap_{i \in I} X \setminus U_i = X \setminus \underbrace{\bigcup_{i \in I} U_i}_{= X} \quad \light \qedhere
	\]
\end{beweis}


\printindex
\printbibliography
\listoffigures
\todototoc
\listoftodos[To-do's und andere Baustellen]
\end{document}
